\documentclass[aspectratio=169, xcolor=dvipsnames]{beamer}

\AtBeginSection[]
{
    \begin{frame}
        \frametitle{Visão Geral}
        \tableofcontents[currentsection]
    \end{frame}
}

\usefonttheme{professionalfonts} % using non standard fonts for beamer
\usefonttheme{serif} % default family is serif
\usepackage{fontspec}
\setmainfont{XCharter}[
    UprightFont    = *-Roman,
    BoldFont       = *-Bold,
    ItalicFont     = *-Italic,
    BoldItalicFont = *-BoldItalic,
]

\usepackage{hyperref}
\usepackage{graphicx} % Allows including images
\usepackage{booktabs} % Allows the use of \toprule, \midrule and \bottomrule in tables
\input{./slides/SimplePlusTheme/beamerthemeSimplePlus.sty}
\input{./slides/SimplePlusTheme/beamerinnerthemeSimplePlus.sty}
\input{./slides/SimplePlusTheme/beamerfontthemeSimplePlus.sty}
\input{./slides/SimplePlusTheme/beamercolorthemeSimplePlus.sty}

\usepackage{listings}
\setbeamercovered{transparent}

\usepackage{biblatex}
\addbibresource{slides/reference.bib}
\renewcommand*{\bibfont}{\footnotesize}
% Tamanho dos slides de "Leituras Recomendadas": maior que a bibliografia
% completa (\bibfont acima, que rende \tiny via a redefinição de \footnotesize)
% para destacar as leituras curadas.
\newcommand{\leiturasize}{\small}

\usepackage[brazilian ]{babel}
\usepackage{csquotes}
\usepackage{pgfplots}
\pgfplotsset{compat=1.18}

\usepackage{emoji}
\usepackage{amsmath}
\usepackage[xcharter,bigdelims,vvarbb]{newtxmath}
% newtxmath's operators font requests the fontspec family `XCharter(0)' in OT1;
% without these substitutions math digits and operator names silently fall back
% to Computer Modern (LaTeX Font Warning: Font shape `OT1/XCharter(0)/m/n' undefined).
\DeclareFontFamily{OT1}{XCharter(0)}{}
\DeclareFontShape{OT1}{XCharter(0)}{m}{n}{<->ssub * XCharter-TLF/m/n}{}
\DeclareFontShape{OT1}{XCharter(0)}{m}{it}{<->ssub * XCharter-TLF/m/it}{}
\DeclareFontShape{OT1}{XCharter(0)}{b}{n}{<->ssub * XCharter-TLF/b/n}{}
\DeclareFontShape{OT1}{XCharter(0)}{bx}{n}{<->ssub * XCharter-TLF/b/n}{}

\usepackage{bm}
\usepackage[table]{xcolor}
\usepackage{xcolor}
\usepackage{algpseudocode}

\usepackage{cancel}
\usepackage{ulem}

\usetikzlibrary{shapes}
\usetikzlibrary{arrows}
\usetikzlibrary {arrows.meta}
\usetikzlibrary{calc,positioning}
\usepgfplotslibrary{fillbetween}
\usetikzlibrary{angles,quotes}
\usetikzlibrary{tikzmark}


\definecolor{PastelRed}{HTML}{FFC1CC}
\definecolor{PastelBlue}{HTML}{C2F0FF}
\definecolor{PastelBlue2}{HTML}{3182BD}

\definecolor{PastelPink}{HTML}{FFCBE1}
\definecolor{PastelGreen}{HTML}{D6E5BD}
\definecolor{PastelYellow}{HTML}{F9E1A8}
\definecolor{PastelBlue}{HTML}{BCD8EC}
\definecolor{PastelPurple}{HTML}{DCCCEC}
\definecolor{PastelOrange}{HTML}{FFDAB4}


\renewcommand{\footnotesize}{\tiny}

\newcommand{\mespor}{
  \ifcase\month
    \or janeiro
    \or fevereiro
    \or março
    \or abril
    \or maio
    \or junho
    \or julho
    \or agosto
    \or setembro
    \or outubro
    \or novembro
    \or dezembro
  \fi
}
% \usepackage{csquotes}
% \usepackage{minted}
% \usemintedstyle{xcode}
\definecolor{vscLightBackground}{HTML}{FFFFFF}
\definecolor{vscLightText}{HTML}{000000}
\definecolor{vscLightGreen}{HTML}{008000}        % For Comments AND Numbers
\definecolor{vscLightRed}{HTML}{A31515}          % For Strings
\definecolor{vscLightTeal}{HTML}{267F99}
\definecolor{vscBlue}{HTML}{0000FF}% For main keywords (if, for, def...)
\definecolor{vscLightPytorch}{HTML}{800080}      % For PyTorch keywords (dark magenta)

\lstdefinestyle{vscode-light}{
    language=Python,
    backgroundcolor=\color{vscLightBackground},
    basicstyle=\ttfamily\scriptsize\color{vscLightText},
    keywordstyle=\color{vscLightTeal},
    commentstyle=\color{vscLightGreen},
    stringstyle=\color{vscLightRed},
    % Style for PyTorch keywords (using keyword group 2)
    keywordstyle=[2]\color{vscLightPytorch}, 
    morekeywords={self, True, False, None},
    morekeywords=[2]{
        torch, Tensor, device, dtype, to, nn, optim, autograd, utils, data,
        Module, Linear, Conv2d, ReLU, CrossEntropyLoss, Adam, SGD,
        Dataset, DataLoader, randn, zeros, ones, arange, cat, stack,
        view, reshape, sum, mean, max, min, matmul, no_grad, backward
    },
    keywordstyle=[3]\color{vscBlue}, 
    morekeywords=[3]{
        model, loss, y, y_hat, x , optimizer
    },
    frame=single,
    tabsize=1,
    showstringspaces=false,
    breaklines=true,
    captionpos=b,
    extendedchars=true,
    numbers=left,
}

\usepackage[cache=true]{minted}
% minted v3+ (TeX Live 2024+) auto-detects the output directory.
% minted v2 (Ubuntu apt TeX Live 2023) needs it set explicitly to
% match latexmkrc's $out_dir = 'build'.
\makeatletter
\@ifpackagelater{minted}{2024/01/01}{}{%
  \def\minted@outputdir{build/}%
}
\makeatother
\setminted{
    fontsize=\scriptsize,
    style=vs,
    breaklines=true,
    numbers=left,
    numbersep=5pt,
    frame=single,
}
\providecommand{\citeauthoryear}[1]{(\citeauthor{#1},~\citeyear{#1})}

\providecommand{\thetav}{\bm{\theta}}

% vetor coluna 2D com colchetes
\providecommand{\cvec}[2]{\begin{bmatrix} #1 \\ #2 \end{bmatrix}}

\providecommand{\varvector}[1]{\mathbf{#1}}
% instance
\providecommand{\X}{\varvector{X}}
\providecommand{\mlinstance}{\varvector{X}}
\providecommand{\mlinstancei}{\varvector{X}^{(i)}}
\providecommand{\mlinstancex}[1]{\varvector{X}^{(#1)}}

% target
\providecommand{\y}{\varvector{y}}
\providecommand{\mltarget}{\varvector{y}}
\providecommand{\mltargeti}{\varvector{y}^{(i)}}
\providecommand{\mltargetx}[1]{\varvector{y}^{(#1)}}

%hats
\providecommand{\yhat}{\hat{\varvector{y}}}
\providecommand{\mltargethat}{\hat{\varvector{y}}}
\providecommand{\mltargethati}{\hat{\varvector{y}}^{(i)}}
\providecommand{\mltargethatx}[1]{\hat{\varvector{y}}^{(#1)}}

% linear reg
\providecommand{\mllinearregression}{\mltargethat=\mlinstance\bm{\theta}}

\providecommand{\tikzarrowcolor}[3]{%
  \begin{tikzpicture}[remember picture, overlay]
    % Arrow body
    \path[fill=#3] ([shift={(#1,#2)}]current page.south west) 
      rectangle ++(0.3, -0.2);
    % Arrow head
    \path[fill=#3] ([shift={(#1+0.3,#2+0.1)}]current page.south west) -- 
         ++(0, -0.4) -- 
         ++(0.2, 0.2) -- 
         cycle;
  \end{tikzpicture}%
}

\providecommand{\tikzarrow}[2]{%
    \tikzarrowcolor{#1}{#2}{Red}
}



\DeclareMathOperator*{\argmax}{argmax}
\DeclareMathOperator*{\argmin}{argmin}



%----------------------------------------------------------------------------------------
%    EXTRA COLORS / COMMANDS FOR THIS LECTURE
%----------------------------------------------------------------------------------------
\definecolor{LightBlue}{HTML}{DBEAFE}
\definecolor{DarkBlue}{HTML}{1565C0}
\definecolor{DarkRed}{HTML}{C62828}
\definecolor{LightGray}{HTML}{F5F5F5}
\definecolor{MidGray}{HTML}{757575}
\definecolor{CellGreen}{HTML}{C8E6C9}
\definecolor{CellYellow}{HTML}{FFF9C4}
\definecolor{CellRed}{HTML}{FFCDD2}
\definecolor{CellBlue}{HTML}{BBDEFB}
\definecolor{CellPurple}{HTML}{E1BEE7}
\definecolor{DarkGreen}{HTML}{1B5E20}

%----------------------------------------------------------------------------------------
%    TITLE PAGE
%----------------------------------------------------------------------------------------

\title{Deep Learning}
\subtitle{\textit{Modelos de Linguagem Decoder-Only}}

\author{Lucas Silveira Kupssinskü}

\institute
{ Machine Learning Theory and Applications Laboratory \\ PUCRS
}
\date{\mespor de \the\year}

%----------------------------------------------------------------------------------------
%    PRESENTATION SLIDES
%----------------------------------------------------------------------------------------

\begin{document}
  \begin{frame}
    \titlepage
  \end{frame}

  \begin{frame}{Visão Geral}
    \tableofcontents
  \end{frame}

  %======================================================================================
  % SECTION 1: REVISÃO DO TRANSFORMER
  %======================================================================================
  \section{Revisão do Transformer}


  \begin{frame}{Arquitetura Original: Encoder-Decoder \cite{vaswani2017attention}}
    \begin{columns}[T]
      \begin{column}{0.50\textwidth}
        \textcolor{DarkGreen}{\textbf{Encoder}} (esquerda):
        \begin{itemize}
          \item Lê toda a entrada \textbf{bidirecional}
          \item Self-attention não mascarada
          \item Produz representações contextuais ricas
        \end{itemize}
        \vspace{0.2em}
        \textcolor{DarkBlue}{\textbf{Decoder}} (direita):
        \begin{itemize}
          \item Gera tokens \textbf{autorregressivamente}
          \item Self-attention \textbf{mascarada} (causal)
          \item \textbf{Cross-attention}: lê saída do encoder
        \end{itemize}
        \vspace{0.3em}
        \begin{alertblock}{Pergunta-chave}
          Precisamos do encoder para geração de linguagem?
        \end{alertblock}
      \end{column}
      \begin{column}{0.46\textwidth}
        \resizebox{\linewidth}{!}{%
        \begin{tikzpicture}[
          enc/.style={draw, rounded corners, minimum width=2.6cm, minimum height=0.65cm,
                      fill=CellGreen, align=center, font=\small},
          dec/.style={draw, rounded corners, minimum width=2.6cm, minimum height=0.65cm,
                      fill=CellYellow, align=center, font=\small},
          arrow/.style={->, >=Stealth, thick},
          node distance=0.5cm
        ]
          \node[enc] (einput) {Input Emb + PE};
          \node[enc, above=of einput] (emha) {Self-Attention};
          \node[enc, above=of emha] (eff) {Feed-Forward};
          \node[above=0.2cm of eff, font=\scriptsize\bfseries] {Encoder $\times N$};

          \draw[arrow] (einput) -- (emha);
          \draw[arrow] (emha) -- (eff);
          \draw[arrow] (eff) -- ++(0,0.5);

          \node[dec, right=1.8cm of einput] (doutput) {Output Emb + PE};
          \node[dec, above=of doutput] (dmask) {Masked Attn};
          \node[dec, above=of dmask] (dcross) {Cross-Attention};
          \node[dec, above=of dcross] (dff) {Feed-Forward};
          \node[above=0.2cm of dff, font=\scriptsize\bfseries] {Decoder $\times N$};
          \node[dec, above=0.6cm of dff] (linear) {Linear + Softmax};

          \draw[arrow] (doutput) -- (dmask);
          \draw[arrow] (dmask) -- (dcross);
          \draw[arrow] (dcross) -- (dff);
          \draw[arrow] (dff) -- (linear);

          \draw[arrow, dashed, DarkRed] (eff.east) ++(0,0) -- ++(0.4,0) |- (dcross.west);
        \end{tikzpicture}}
      \end{column}
    \end{columns}
  \end{frame}

  \begin{frame}{Multi-Head Self-Attention}
    \begin{columns}[T]
      \begin{column}{0.50\textwidth}
        Dado um token, computamos queries, keys e values:
        \begin{block}{Atenção escalada}
          $\text{Attn}(Q,K,V) = \text{softmax}\!\left(\dfrac{QK^\top}{\sqrt{d_k}}\right)V$
        \end{block}
        \vspace{0.3em}
        Cada cabeça aprende padrões diferentes:
        \[
          \text{MHA} = \text{Concat}(\text{head}_1,\ldots,\text{head}_H)\,W^O
        \]
      \end{column}
      \begin{column}{0.46\textwidth}
        \begin{itemize}
          \item $H$ cabeças $\times$ $d_k = d_\text{model}/H$
          \item Custo: $\mathcal{O}(n^2 d)$ por camada
          \item Paralelizável sobre toda a sequência
        \end{itemize}
        \vspace{0.4em}
        \begin{exampleblock}{Máscara causal (decoder)}
          Posição $i$ só pode atender a $j \le i$:
          \[
            M_{ij} = \begin{cases} 0 & j \le i \\ -\infty & j > i \end{cases}
          \]
        \end{exampleblock}
      \end{column}
    \end{columns}
  \end{frame}

  \begin{frame}{Bloco Transformer: FFN, Resíduo e Normalização}
    \begin{columns}[T]
      \begin{column}{0.50\textwidth}
        Cada camada contém:
        \begin{enumerate}
          \item \textbf{(Masked) Self-Attention} + resíduo + norma
          \item \textbf{Feed-Forward} + resíduo + norma
        \end{enumerate}
        \vspace{0.3em}
        \begin{block}{FFN de dois estágios}
          $\text{FFN}(x) = \max(0,\, xW_1 + b_1)\,W_2 + b_2$
        \end{block}
        Dimensão interna: $d_\text{ff} \approx 4\,d_\text{model}$\\
        {\footnotesize (Valor empírico de \cite{vaswani2017attention}; com SwiGLU parte-se de $\tfrac{8}{3}\,d_\text{model}$, arredondado para múltiplos convenientes.)}
      \end{column}
      \begin{column}{0.46\textwidth}
        \begin{block}{Add \& Norm (por subcamada)}
          $x \leftarrow \text{Norm}\bigl(x + \text{SubLayer}(x)\bigr)$\\[2pt]
          O resíduo preserva o sinal; a \textbf{norma} estabiliza as ativações.
        \end{block}
        \vspace{0.3em}
        Empilhar $N$ desses blocos $=$ Transformer.\\[4pt]
        {\footnotesize Os tipos de norma (Layer/RMS Norm) e sua posição (Pre/Post-Norm) vêm na próxima seção.}
      \end{column}
    \end{columns}
  \end{frame}

  \begin{frame}{FFN Padrão: Exemplo Numérico (ReLU)}
    \begin{columns}[T]
      % ── LEFT: single-path computation flow ──────────────────────────────
      \begin{column}{0.52\textwidth}
        \centering
        \resizebox{0.78\linewidth}{!}{%
        \begin{tikzpicture}[
          vec/.style={draw, rounded corners, minimum width=3.0cm, minimum height=0.52cm,
                      fill=CellBlue, align=center, font=\small},
          lin/.style={draw, rounded corners, minimum width=2.6cm, minimum height=0.52cm,
                      fill=CellYellow, align=center, font=\small},
          act/.style={draw, rounded corners, minimum width=2.6cm, minimum height=0.52cm,
                      fill=CellRed, align=center, font=\small},
          arr/.style={->, >=Stealth, thick},
          ann/.style={font=\scriptsize, gray},
        ]
          \node[vec] (x)    at (0, 5.2) {$x \in \mathbb{R}^{d_\mathrm{model}}$};
          \node[lin] (w1)   at (0, 3.9) {$x W_1$};
          \node[act] (relu) at (0, 2.6) {$\mathrm{ReLU}(x W_1)$};
          \node[lin] (w2)   at (0, 1.3) {$(\cdots)\,W_2$};
          \node[vec] (y)    at (0, 0.0) {$y \in \mathbb{R}^{d_\mathrm{model}}$};

          \draw[arr] (x)    -- (w1);
          \draw[arr] (w1)   -- (relu);
          \draw[arr] (relu) -- (w2);
          \draw[arr] (w2)   -- (y);

          \node[ann, right=0.12cm of w1]   {$d_m \to d_{ff}{=}4d_m$};
          \node[ann, right=0.12cm of relu] {$\max(0,\,z)$};
          \node[ann, right=0.12cm of w2]   {$d_{ff} \to d_m$};
        \end{tikzpicture}}
      \end{column}

      % ── RIGHT: numerical example ─────────────────────────────────────
      \begin{column}{0.44\textwidth}
        \footnotesize\textbf{Exemplo numérico}\\[0.1em]
        \scriptsize$d_\mathrm{model}=3,\quad d_{ff}=4\times3=12$\\[0.5em]

        {\setlength{\fboxsep}{3pt}%
        \colorbox{CellBlue}{$x = [\,1.0,\;0.5,\;{-0.5}\,]$}\\[0.35em]
        \colorbox{CellYellow}{$xW_1 = [\,2.0,\;{-1.0},\;0.5,\;\ldots\,]$\scriptsize\textit{ (dim 12)}}\\[0.35em]
        \colorbox{CellRed}{$\mathrm{ReLU}(xW_1) = [\,2.0,\;\mathbf{0.0},\;0.5,\;\ldots\,]$}\\[0.3em]
        }%
        \textcolor{red!70!black}{\tiny$\Rightarrow$ ReLU \textbf{zera} o valor negativo $-1.0$!}\\[0.5em]

        {\setlength{\fboxsep}{3pt}%
        \colorbox{CellYellow}{$[\,\cdots\,]W_2\;\approx\;y = [\,0.7,\;{-0.2},\;0.4\,]$}
        }
      \end{column}
    \end{columns}
  \end{frame}

  \begin{frame}{SwiGLU: A FFN com \textit{Gate} das LLMs mais atuais}
    \begin{columns}[T]
      \begin{column}{0.50\textwidth}
        FFN clássica usa \textbf{ReLU}:
        \begin{block}{ReLU-FFN (Vaswani 2017)}
          $\text{FFN}(x) = \max(0,\, xW_1)\,W_2$
        \end{block}
        \vspace{0.4em}
        Llama 2/3, PaLM, Mistral usam \textbf{SwiGLU} \cite{shazeer2020swiglu}:
        \begin{block}{SwiGLU-FFN}
          $\text{FFN}(x) = \bigl(\text{SiLU}(xW_1) \odot xW_3\bigr)\,W_2$
        \end{block}
        onde $\text{SiLU}(z) = z \cdot \sigma(z)$ (\textit{Swish}).\\[0.3em]
        O \textbf{portão} $xW_3$ filtra dinamicamente a informação.
      \end{column}
      \begin{column}{0.46\textwidth}
        \begin{exampleblock}{Por que $\tfrac{8}{3}\,d_\text{model}$?}
          SwiGLU usa \textbf{3 matrizes} em vez de 2.\\
          Para manter o mesmo número de parâmetros:\\[0.2em]
          $3 \times d_\text{ff} \cdot d = 2 \times 4d \cdot d$\\[0.2em]
          $\Rightarrow d_\text{ff} = \tfrac{8}{3}\,d_\text{model}$
        \end{exampleblock}
        \vspace{0.3em}
        \begin{itemize}
          \item SwiGLU supera ReLU e GELU empiricamente
          \item Motivação teórica: portão aprende \textit{quais} dimensões ativar
          \item Alternativas: GeGLU, ReGLU
        \end{itemize}
      \end{column}
    \end{columns}
  \end{frame}

  \begin{frame}{SwiGLU: Diagrama Computacional e Dimensões}
    \begin{columns}[T]
      % ── LEFT: computation flow ──────────────────────────────────────────
      \begin{column}{0.55\textwidth}
        \centering
        \resizebox{0.88\linewidth}{!}{%
        \begin{tikzpicture}[
          % Boxes for tensors (input/output)
          vec/.style={draw, rounded corners, minimum width=3.0cm, minimum height=0.52cm,
                      fill=CellBlue, align=center, font=\small},
          % Boxes for linear projections
          lin/.style={draw, rounded corners, minimum width=2.4cm, minimum height=0.52cm,
                      fill=CellYellow, align=center, font=\small},
          % Gate branch projection
          gate/.style={draw, rounded corners, minimum width=2.6cm, minimum height=0.52cm,
                       fill=CellGreen, align=center, font=\small},
          % SiLU activation
          act/.style={draw, rounded corners, minimum width=2.4cm, minimum height=0.52cm,
                      fill=CellRed, align=center, font=\small},
          % Element-wise multiply node
          mul/.style={draw, circle, minimum size=0.52cm, fill=PastelOrange,
                      inner sep=0pt, font=\normalsize, thick},
          % Arrows
          arr/.style={->, >=Stealth, thick},
          % Small grey annotation labels
          ann/.style={font=\scriptsize, gray},
        ]
          % ── NODES ──────────────────────────────────────────────────────
          % Input token representation
          \node[vec]  (x)    at (1.6, 6.2) {$x \in \mathbb{R}^{d_\mathrm{model}}$};

          % Left branch: compute path via W1 (will be activated)
          \node[lin]  (w1)   at (0.0, 4.9) {$x W_1$};
          % Right branch: gate path via W3 — no activation
          \node[gate] (w3)   at (3.2, 4.9) {$x W_3$\;\textit{(portão)}};

          % SiLU activation applied only to the compute path
          \node[act]  (silu) at (0.0, 3.6) {$\mathrm{SiLU}(x W_1)$};

          % Element-wise product: gating = compute ⊙ gate
          \node[mul]  (mul)  at (1.6, 2.4) {$\odot$};

          % Final linear projection back to d_model
          \node[lin]  (w2)   at (1.6, 1.2) {$(\cdots)\,W_2$};

          % Output token representation
          \node[vec]  (y)    at (1.6, 0.0) {$y \in \mathbb{R}^{d_\mathrm{model}}$};

          % ── ARROWS ─────────────────────────────────────────────────────
          % x splits into two parallel branches
          \draw[arr] (x.south) -- ++(0,-0.42) -| (w1.north);
          \draw[arr] (x.south) -- ++(0,-0.42) -| (w3.north);

          % Compute path: W1 → SiLU
          \draw[arr] (w1) -- (silu);

          % SiLU output enters element-wise multiply from the left
          \draw[arr] (silu.south) -- ++(0,-0.35) -| (mul.west);

          % Gate path: W3 descends and enters ⊙ from the right
          \draw[arr] (w3.south) -- ++(0,-1.1) -| (mul.east);

          % Merged output passes through W2 to produce y
          \draw[arr] (mul) -- (w2);
          \draw[arr] (w2)  -- (y);

          % ── DIMENSION ANNOTATIONS ─────────────────────────────────────
          % Each matrix shape shown in grey next to its node
          \node[ann, left=0.08cm of w1]  {$d_m \to d_{ff}$};
          \node[ann, right=0.08cm of w3] {$d_m \to d_{ff}$};
          \node[ann, right=0.08cm of w2] {$d_{ff} \to d_m$};

        \end{tikzpicture}}
      \end{column}

      % ── RIGHT: numerical example ─────────────────────────────────────
      \begin{column}{0.41\textwidth}
        \footnotesize\textbf{Exemplo numérico}\\[0.1em]
        \scriptsize$d_\mathrm{model}=3,\quad d_{ff}=\tfrac{8}{3}\!\times\!3=8$\\[0.5em]

        {\setlength{\fboxsep}{3pt}%
        \colorbox{CellBlue}{$x = [\,1.0,\;0.5,\;{-0.5}\,]$}\\[0.35em]
        \colorbox{CellYellow}{$xW_1 = [\,2.0,\;{-1.0},\;0.5,\;1.2,\;\ldots\,]$}\\[0.15em]
        \colorbox{CellRed}{$\mathrm{SiLU}(xW_1) \approx [\,1.76,\;{-0.27},\;0.31,\;0.92,\;\ldots\,]$}\\[0.45em]
        \colorbox{CellGreen}{$xW_3 = [\,0.8,\;\mathbf{0.0},\;1.2,\;0.5,\;\ldots\,]$\;\textit{(portão)}}\\[0.6em]
        }%

        \colorbox{PastelOrange}{%
          \begin{minipage}{0.88\linewidth}
            \scriptsize
            $[\,1.76,\;{-0.27},\;0.31,\;0.92,\;\ldots\,]$\\
            $\odot\;[\,0.8,\;\mathbf{0.0},\;1.2,\;0.5,\;\ldots\,]$\\[0.15em]
            $=[\,1.41,\;\mathbf{0.00},\;0.37,\;0.46,\;\ldots\,]$\\[0.2em]
            \textcolor{red!70!black}{\tiny$\Rightarrow$ portão \textbf{0.0} suprime a dim.\,2!}
          \end{minipage}%
        }\\[0.6em]

        {\setlength{\fboxsep}{3pt}%
        \colorbox{CellYellow}{$[\,\cdots\,]W_2\;\approx\;y = [\,0.7,\;{-0.2},\;0.4\,]$}
        }
      \end{column}
    \end{columns}
  \end{frame}

  %======================================================================================
  % SECTION 4: LAYER NORM E RMS NORM
  %======================================================================================
  \section{Normalização: Layer Norm e RMS Norm}

  \begin{frame}{Por que Normalizar as Ativações?}
    \begin{columns}[T]
      \begin{column}{0.50\textwidth}
        Durante o treino de redes profundas:
        \begin{itemize}
          \item Distribuição das ativações muda a cada camada (\textit{internal covariate shift})
          \item Gradientes explodem ou desaparecem
          \item Taxas de aprendizado devem ser muito pequenas
        \end{itemize}
        \vspace{0.4em}
        \textbf{Batch Norm}: normaliza sobre o batch, problemático para sequências de tamanho variável.
      \end{column}
      \begin{column}{0.46\textwidth}
        \begin{block}{Solução para Transformers}
          \textbf{Layer Normalization} \cite{ba2016layernorm}: normaliza sobre as \textit{features} de cada token, sem dependência do batch.
        \end{block}
        \vspace{0.3em}
        Permite treino estável de Transformers muito profundos com learning rates maiores.
      \end{column}
    \end{columns}
  \end{frame}

  \begin{frame}{Layer Normalization \cite{ba2016layernorm}}
    \begin{columns}[T]
      \begin{column}{0.52\textwidth}
        Dado o vetor de ativações $\mathbf{x} \in \mathbb{R}^d$ de um token:
        \[
          \mu = \frac{1}{d}\sum_{i=1}^d x_i, \qquad
          \sigma^2 = \frac{1}{d}\sum_{i=1}^d (x_i - \mu)^2
        \]
        \[
          \hat{x}_i = \frac{x_i - \mu}{\sqrt{\sigma^2 + \epsilon}}
        \]
        \[
          \text{LN}(\mathbf{x}) = \boldsymbol{\gamma} \odot \hat{\mathbf{x}} + \boldsymbol{\beta}
        \]
        $\boldsymbol{\gamma}, \boldsymbol{\beta} \in \mathbb{R}^d$: parâmetros aprendidos (\textit{gain} e \textit{bias}).
      \end{column}
      \begin{column}{0.44\textwidth}
        \begin{itemize}
          \item Independente do tamanho do batch
          \item Funciona com sequências de tamanho variável
          \item Aplicada \textbf{por token} (ao longo da dimensão $d$)
        \end{itemize}
        \vspace{0.3em}
        \begin{exampleblock}{Nota}
          A normalização ocorre ao longo da dimensão de features $d$ de cada token, não ao longo da sequência.
        \end{exampleblock}
      \end{column}
    \end{columns}
  \end{frame}

  \begin{frame}{Layer Norm, RMS Norm e Batch Norm: Exemplo Numérico}
    Lote de 2 tokens, cada um com $d=4$ \textit{features}:
    \begin{center}
      $X=\begin{bmatrix}1&3&5&7\\9&7&5&3\end{bmatrix}$\qquad
      {\small\textcolor{DarkBlue}{$\rightarrow$ LN/RMS: por linha (token)}\quad\textcolor{DarkRed}{$\downarrow$ BN: por coluna (feature)}}
    \end{center}
    \begin{columns}[T]
      \begin{column}{0.52\textwidth}
        \begin{block}{LayerNorm --- token 1 (eixo $d$)}
          $\mu=4,\quad\sigma^2=5$\\[2pt]
          $\hat{\mathbf{x}}=\dfrac{\mathbf{x}-\mu}{\sqrt{\sigma^2}}=[-1.34,\,-0.45,\,0.45,\,1.34]$\\[2pt]
          {\footnotesize média $0$ e variância $1$, usando só as features do token}
        \end{block}
        \begin{block}{RMSNorm --- token 1 (eixo $d$)}
          $\text{RMS}=\sqrt{\tfrac{1}{4}\textstyle\sum_i x_i^2}=4.58$\\[2pt]
          $\mathbf{x}/\text{RMS}=[0.22,\,0.66,\,1.09,\,1.53]$\ {\footnotesize(só reescala)}
        \end{block}
      \end{column}
      \begin{column}{0.44\textwidth}
        \begin{alertblock}{BatchNorm --- feature 1 (lote)}
          coluna $1=[1,\,9]$:\ $\mu=5,\ \sigma^2=16$\\[2pt]
          $\hat{\mathbf{x}}=[-1,\,1]$\\[2pt]
          {\footnotesize usa estatísticas \textbf{do lote}, depende do batch}
        \end{alertblock}
        \begin{exampleblock}{Conclusão}
          LN/RMS normalizam \textbf{por token}; BN \textbf{por feature}. Transformers usam LN/RMS por independerem do batch.
        \end{exampleblock}
      \end{column}
    \end{columns}
  \end{frame}


  \begin{frame}{Pre-Norm vs.\ Post-Norm}
    \begin{columns}[T]
      \begin{column}{0.48\textwidth}
        \begin{block}{Post-Norm (Vaswani 2017)}
          $x \leftarrow \text{LN}(x + \text{SubLayer}(x))$\\[0.3em]
          Normalização \textbf{depois} da operação.\\
          \textcolor{DarkRed}{$\to$ gradientes instáveis com $N$ grande; requer \textit{warm-up} lento.}
        \end{block}
      \end{column}
      \begin{column}{0.48\textwidth}
        \begin{block}{Pre-Norm (GPT-2 em diante)}
          $x \leftarrow x + \text{SubLayer}(\text{LN}(x))$\\[0.3em]
          Normalização \textbf{antes} da operação.\\
          \textcolor{DarkGreen}{$\to$ treino mais estável; gradientes fluem melhor pela conexão residual.}
        \end{block}
      \end{column}
    \end{columns}
    \vspace{0.5em}
    \begin{center}
      \textbf{Hoje}, quase todos os LLMs usam \textbf{Pre-Norm}.\\
      A conexão residual ``bypassa'' a normalização, preservando o sinal ao longo das camadas.
    \end{center}
  \end{frame}

  \begin{frame}{RMS Norm: Simplificação do Layer Norm \cite{zhang2019rmsnorm}}
    \begin{columns}[T]
      \begin{column}{0.52\textwidth}
        \textbf{Zhang \& Sennrich (2019)}: o re-centramento (subtração de $\mu$) é desnecessário quando a invariância de escala é suficiente.\\[0.4em]
        Normaliza apenas pela raiz do valor quadrático médio:
        \[
          \text{RMS}(\mathbf{x}) = \sqrt{\frac{1}{d}\sum_{i=1}^d x_i^2}
        \]
        \[
          \text{RMSNorm}(\mathbf{x}) = \frac{\mathbf{x}}{\text{RMS}(\mathbf{x})} \odot \boldsymbol{\gamma}
        \]
        \textbf{Sem} $\mu$, \textbf{sem} $\boldsymbol{\beta}$, mais simples e rápido.
      \end{column}
      \begin{column}{0.44\textwidth}
        \begin{itemize}
          \item Mais rápido: sem cálculo de média
          \item Mesma qualidade empírica que LayerNorm
          \item Menos parâmetros (sem $\boldsymbol{\beta}$)
        \end{itemize}
        \vspace{0.3em}
        \begin{exampleblock}{Adotado por}
          Llama 1/2/3, Mistral, Mixtral, Gemma, Falcon, PaLM, \ldots
        \end{exampleblock}
        \vspace{0.2em}
        Citado em \cite{stanfordcme295} como padrão atual para LLMs eficientes.
      \end{column}
    \end{columns}
  \end{frame}

  \begin{frame}{Qual Normalização Cada Modelo Usa?}
    \begin{center}
    \resizebox{0.88\linewidth}{!}{%
    \begin{tabular}{lllll}
      \toprule
      \textbf{Modelo} & \textbf{Norma} & \textbf{Posição} & \textbf{Bias $\boldsymbol{\beta}$?} & \textbf{Observação} \\
      \midrule
      BERT            & LayerNorm & Post-Norm & Sim & Encoder bidirecional \\
      GPT-2           & LayerNorm & Pre-Norm  & Sim & Decoder-only \\
      GPT-3           & LayerNorm & Pre-Norm  & Sim & Decoder-only \\
      Llama 1/2/3     & RMSNorm   & Pre-Norm  & Não & Decoder-only \\
      Mistral 7B      & RMSNorm   & Pre-Norm  & Não & Decoder-only \\
      Gemma           & RMSNorm   & Pre+Post  & Não & ``Sandwich'' norm, estabiliza treino profundo \\
      PaLM            & RMSNorm   & Pre-Norm  & Não & Decoder-only \\
      \bottomrule
    \end{tabular}}
    \end{center}
    \vspace{0.3em}
    \begin{center}
      \textcolor{DarkBlue}{\textbf{Tendência: RMSNorm + Pre-Norm para LLMs modernos.}}
    \end{center}
  \end{frame}

  %======================================================================================
  % SECTION: DO ENCODER-DECODER AO DECODER-ONLY
  %======================================================================================
  \section{Do Encoder-Decoder ao Decoder-Only}

  \begin{frame}{O que Encoder e Decoder fazem?}
    \begin{columns}[T]
      \begin{column}{0.48\textwidth}
        \begin{block}{Encoder: compreensão}
          \begin{itemize}
            \item Processa \textbf{toda a entrada} de uma vez
            \item Atenção bidirecional
            \item Produz vetor de contexto rico
            \item Exemplos: BERT, RoBERTa
          \end{itemize}
        \end{block}
      \end{column}
      \begin{column}{0.48\textwidth}
        \begin{block}{Decoder: geração}
          \begin{itemize}
            \item Gera \textbf{um token por vez}
            \item Atenção causal (só vê o passado)
            \item Recebe contexto via cross-attention
            \item Exemplos: GPT, T5-decoder
          \end{itemize}
        \end{block}
      \end{column}
    \end{columns}
    \vspace{0.5em}
    \begin{center}
      \textcolor{DarkRed}{\textbf{Para LLMs, o encoder é desnecessário.}}\\
      \vspace{0.2em}
      O decoder processa o \textit{prompt} como parte da sequência.
    \end{center}
  \end{frame}

  \begin{frame}{Máscara Causal vs.\ Atenção Bidirecional}
    \begin{columns}[T]
      \begin{column}{0.50\textwidth}
        \textbf{Bidirecional (BERT/Encoder)}:\\
        cada posição vê \textit{toda} a sequência.
        \vspace{0.3em}

        \resizebox{0.85\linewidth}{!}{%
        \begin{tabular}{c|cccc}
          & $t_1$ & $t_2$ & $t_3$ & $t_4$ \\\hline
          $t_1$ & \cellcolor{CellGreen}$\checkmark$ & \cellcolor{CellGreen}$\checkmark$ & \cellcolor{CellGreen}$\checkmark$ & \cellcolor{CellGreen}$\checkmark$ \\
          $t_2$ & \cellcolor{CellGreen}$\checkmark$ & \cellcolor{CellGreen}$\checkmark$ & \cellcolor{CellGreen}$\checkmark$ & \cellcolor{CellGreen}$\checkmark$ \\
          $t_3$ & \cellcolor{CellGreen}$\checkmark$ & \cellcolor{CellGreen}$\checkmark$ & \cellcolor{CellGreen}$\checkmark$ & \cellcolor{CellGreen}$\checkmark$ \\
          $t_4$ & \cellcolor{CellGreen}$\checkmark$ & \cellcolor{CellGreen}$\checkmark$ & \cellcolor{CellGreen}$\checkmark$ & \cellcolor{CellGreen}$\checkmark$ \\
        \end{tabular}}
      \end{column}
      \begin{column}{0.46\textwidth}
        \textbf{Causal (GPT/Decoder-Only)}:\\
        posição $i$ só vê $j \le i$.
        \vspace{0.3em}

        \resizebox{0.85\linewidth}{!}{%
        \begin{tabular}{c|cccc}
          & $t_1$ & $t_2$ & $t_3$ & $t_4$ \\\hline
          $t_1$ & \cellcolor{CellGreen}$\checkmark$ & \cellcolor{CellRed}$\times$ & \cellcolor{CellRed}$\times$ & \cellcolor{CellRed}$\times$ \\
          $t_2$ & \cellcolor{CellGreen}$\checkmark$ & \cellcolor{CellGreen}$\checkmark$ & \cellcolor{CellRed}$\times$ & \cellcolor{CellRed}$\times$ \\
          $t_3$ & \cellcolor{CellGreen}$\checkmark$ & \cellcolor{CellGreen}$\checkmark$ & \cellcolor{CellGreen}$\checkmark$ & \cellcolor{CellRed}$\times$ \\
          $t_4$ & \cellcolor{CellGreen}$\checkmark$ & \cellcolor{CellGreen}$\checkmark$ & \cellcolor{CellGreen}$\checkmark$ & \cellcolor{CellGreen}$\checkmark$ \\
        \end{tabular}}
        \vspace{0.3em}
        Garante que $t_4$ não ``veja o futuro'' durante o treino.
      \end{column}
    \end{columns}
  \end{frame}

  \begin{frame}{Do Encoder-Decoder para o Decoder-Only}
    \begin{columns}[T]
      \begin{column}{0.50\textwidth}
        \textbf{Ideia central}: treine apenas o decoder.\\[0.5em]
        O modelo recebe o texto de entrada como \textit{prefixo} e continua gerando autorregressivamente.\\[0.5em]
        \begin{itemize}
          \item Sem cross-attention
          \item Sem encoder separado
          \item Toda informação está nas posições anteriores
          \item Arquitetura mais simples
        \end{itemize}
      \end{column}
      \begin{column}{0.46\textwidth}
        \centering
        \begin{tikzpicture}[
          block/.style={draw, rounded corners, minimum width=4.5cm, minimum height=0.50cm,
                        fill=CellYellow, align=center, font=\footnotesize},
          norm/.style={draw, rounded corners, minimum width=4.5cm, minimum height=0.50cm,
                       fill=CellBlue, align=center, font=\footnotesize},
          add/.style={draw, circle, minimum size=0.32cm, fill=white,
                      inner sep=0pt, font=\footnotesize},
          arrow/.style={->, >=Stealth, thick},
          node distance=0.28cm
        ]
          \node[norm]  (n1)   {RMSNorm};
          \node[block, above=of n1]   (attn) {Masked Self-Attention};
          \node[add,   above=of attn] (add1) {$+$};
          \node[norm,  above=of add1] (n2)   {RMSNorm};
          \node[block, above=of n2]   (ff)   {Feed-Forward (MLP)};
          \node[add,   above=of ff]   (add2) {$+$};
          \node[block, above=of add2] (lm)   {LM Head (Linear)};
          \node[below=0.4cm of n1, font=\footnotesize] (inp) {tokens de entrada};

          % Main forward flow
          \draw[arrow] (inp)  -- (n1);
          \draw[arrow] (n1)   -- (attn);
          \draw[arrow] (attn) -- (add1);
          \draw[arrow] (add1) -- (n2);
          \draw[arrow] (n2)   -- (ff);
          \draw[arrow] (ff)   -- (add2);
          \draw[arrow] (add2) -- (lm);

          % Reference point just outside the left edge of all blocks
          \coordinate (resLeft) at ([xshift=-0.4cm]n1.west);

          % Residual 1: inp → left of blocks → add1
          \draw[->, >=Stealth, thick, DarkRed]
            (inp.west) -- (inp.west -| resLeft) -- (add1.west -| resLeft) -- (add1.west);
          % Residual 2: add1 → left of blocks → add2
          \draw[->, >=Stealth, thick, DarkRed]
            (add1.west) -- (add1.west -| resLeft) -- (add2.west -| resLeft) -- (add2.west);
          \node[font=\tiny, DarkRed, left] at ($(add1.west -| resLeft)!0.5!(add2.west -| resLeft)$) {resíduo};
        \end{tikzpicture}
        \vspace{0.2em}
        \footnotesize Repetido $N$ vezes (bloco).
      \end{column}
    \end{columns}
  \end{frame}

  %======================================================================================
  % SECTION 2: A FAMÍLIA GPT
  %======================================================================================
  \section{A Família GPT e o Paradigma Decoder-Only}

  \begin{frame}{GPT-1: Pré-Treino + Fine-Tuning \cite{radford2018gpt1}}
    \begin{columns}[T]
      \begin{column}{0.52\textwidth}
        \textbf{Radford et al.\ (2018)}, OpenAI\\[0.4em]
        \textbf{Objetivo de pré-treino}: modelagem de linguagem causal
        \[
          \mathcal{L} = -\sum_{t} \log P(x_t \mid x_1,\ldots,x_{t-1};\,\Theta)
        \]
        Depois: \textit{fine-tuning} supervisionado para cada tarefa.
      \end{column}
      \begin{column}{0.44\textwidth}
        \begin{itemize}
          \item 117M parâmetros
          \item 12 camadas, $d=768$, 12 cabeças
          \item BooksCorpus ($\approx$800M tokens)
          \item Tokenização BPE
        \end{itemize}
        \vspace{0.3em}
        \begin{exampleblock}{Contribuição}
          Mostrou que pré-treino auto-supervisionado em texto melhora fortemente o fine-tuning.
        \end{exampleblock}
      \end{column}
    \end{columns}
  \end{frame}

  \begin{frame}{GPT-2: Geração Zero-Shot \cite{radford2019gpt2}}
    \begin{columns}[T]
      \begin{column}{0.52\textwidth}
        \textbf{Radford et al.\ (2019)}\\[0.4em]
        \textbf{Hipótese}: um LM suficientemente grande aprende a realizar tarefas \textbf{sem fine-tuning}.\\[0.4em]
        \textit{``Language models are unsupervised multitask learners''}\\[0.4em]
        Novidades técnicas:
        \begin{itemize}
          \item Pre-Norm ao invés de Post-Norm
          \item Contexto: 1024 tokens
          \item Vocabulário: 50257 tokens (BPE\footnote{\scriptsize BPE = \textit{Byte-Pair Encoding}: segmenta palavras em sub-palavras frequentes, reduzindo o tamanho do vocabulário.})
        \end{itemize}
      \end{column}
      \begin{column}{0.44\textwidth}
        \resizebox{\linewidth}{!}{%
        \begin{tabular}{lrrrr}
          \toprule
          \textbf{Versão} & \textbf{Params} & \textbf{Camadas} & $d$ & $H$ \\
          \midrule
          Small  & 117M  & 12 & 768  & 12 \\
          Medium & 345M  & 24 & 1024 & 16 \\
          Large  & 762M  & 36 & 1280 & 20 \\
          XL     & 1.5B  & 48 & 1600 & 25 \\
          \bottomrule
        \end{tabular}}
        \vspace{0.3em}
        Dataset: WebText (40GB de texto web filtrado).
      \end{column}
    \end{columns}
  \end{frame}

  \begin{frame}{GPT-3: Zero e Few-shot learning\cite{brown2020gpt3}}
    \begin{columns}[T]
      \begin{column}{0.52\textwidth}
        \textbf{Brown et al.\ (2020)}, 175 bilhões de parâmetros\\[0.4em]
        \textbf{Few-shot prompting}: exemplos fornecidos \textit{no contexto}, sem atualização de pesos.
        \vspace{0.3em}
        \begin{block}{Modos de uso em contexto}
          \begin{itemize}
            \item \textbf{Zero-shot}: só a instrução
            \item \textbf{One-shot}: 1 exemplo no prompt
            \item \textbf{Few-shot}: K exemplos no prompt
          \end{itemize}
        \end{block}
      \end{column}
      \begin{column}{0.44\textwidth}
        \begin{itemize}
          \item 96 camadas, $d=12288$, 96 cabeças
          \item Contexto: 2048 tokens
          \item Treinado em 300B tokens
          \item \textbf{Capacidades emergentes}: habilidades que \textit{surgem abruptamente} com escala (aritmética, tradução, raciocínio)
        \end{itemize}
        \vspace{0.2em}
        \begin{alertblock}{Mudança de paradigma}
          A partir do GPT-3, o fine-tuning por tarefa foi substituído pela engenharia de \textit{prompts}.
        \end{alertblock}
      \end{column}
    \end{columns}
  \end{frame}

  \begin{frame}{Bloco Transformer Decoder-Only em Detalhe}
    \begin{columns}[T]
      \begin{column}{0.50\textwidth}
        O bloco básico repete-se $N$ vezes:\\[0.3em]
        \begin{enumerate}
          \item \textbf{Normalização} (Pre-Norm)
          \item \textbf{Masked Self-Attention}
          \item Conexão residual: $x \mathrel{+}= \text{Attn}(\cdot)$
          \item \textbf{Normalização}
          \item \textbf{Feed-Forward} (MLP)
          \item Conexão residual: $x \mathrel{+}= \text{FFN}(\cdot)$
        \end{enumerate}
        \vspace{0.3em}
        \textbf{Sem} cross-attention. Toda informação está nas $n$ posições anteriores da sequência.
      \end{column}
      \begin{column}{0.46\textwidth}
        \begin{exampleblock}{KV-Cache}
          Durante inferência, K e V já calculados são armazenados.\\[0.2em]
          \textbf{Sem cache}: custo $\mathcal{O}(n^2 d)$ por token\\
          \textbf{Com cache}: custo $\mathcal{O}(n d)$ por token\\[0.2em]
          Essencial para geração eficiente em sequências longas.
        \end{exampleblock}
      \end{column}
    \end{columns}
  \end{frame}

  \begin{frame}{KV-Cache: Como Funciona na Prática}
    \begin{columns}[T]
      \begin{column}{0.50\textwidth}
        Gerando a sequência \textit{``O rato roeu a rolha''}:
        \vspace{0.3em}
        \begin{enumerate}
          \item \textbf{Prefill}: alimenta o \textit{prompt} inteiro de uma vez\\
                $\to$ computa e \textit{armazena} K, V de cada token
          \item \textbf{Decode passo 1}: prediz próximo token\\
                $\to$ computa K, V \textit{só} do token novo\\
                $\to$ reutiliza K, V dos tokens anteriores do cache
          \item \textbf{Decode passo 2, 3, \ldots}: idem\\
                $\to$ cache cresce a cada passo gerado
        \end{enumerate}
      \end{column}
      \begin{column}{0.46\textwidth}
        \begin{block}{Sem cache (naïve)}
          Para o $n$-ésimo token:\\
          Recalcula K, V de todos os $n$ tokens\\
          Custo total: $\mathcal{O}(n^2 d)$
        \end{block}
        \vspace{0.3em}
        \begin{block}{Com cache}
          Computa K, V apenas do token novo\\
          Concatena com o cache existente\\
          Custo por token: $\mathcal{O}(n d)$
        \end{block}
        \vspace{0.2em}
        {\footnotesize Memória do cache $\propto n \cdot L \cdot d_k$, cresce com a sequência e número de camadas.}
      \end{column}
    \end{columns}
  \end{frame}

  % ── KV-Cache step-by-step TikZ visualisation ──────────────────────────────

  \begin{frame}{KV-Cache: Visualização}{Prefill}
    \small Exemplo: \textit{``O rato roeu a rolha''}, todos os K,V do prompt são computados de uma vez.\\[0.3em]
    \begin{center}
    \begin{tikzpicture}[
      tok/.style={draw, rounded corners, minimum width=1.6cm, minimum height=0.52cm,
                  align=center, font=\small},
      kv/.style={draw, minimum width=1.6cm, minimum height=0.40cm,
                 align=center, font=\scriptsize},
      arr/.style={->, >=Stealth, thick},
    ]
      % ── Prompt tokens ─────────────────────────────────────────────────
      \node[tok, fill=CellBlue] (t1) at (0.8, 3.5) {\texttt{O}};
      \node[tok, fill=CellBlue] (t2) at (2.9, 3.5) {\texttt{rato}};
      \node[tok, fill=CellBlue] (t3) at (5.0, 3.5) {\texttt{roeu}};
      \node[font=\scriptsize, DarkBlue, above=0.05cm of t2] {\textbf{Prompt} (3 tokens)};

      % ── K row (being computed = CellYellow) ───────────────────────────
      \node[kv, fill=CellYellow] (k1) at (0.8, 2.45) {$K_1$};
      \node[kv, fill=CellYellow] (k2) at (2.9, 2.45) {$K_2$};
      \node[kv, fill=CellYellow] (k3) at (5.0, 2.45) {$K_3$};

      % ── V row (being computed = CellYellow) ───────────────────────────
      \node[kv, fill=CellYellow] (v1) at (0.8, 1.75) {$V_1$};
      \node[kv, fill=CellYellow] (v2) at (2.9, 1.75) {$V_2$};
      \node[kv, fill=CellYellow] (v3) at (5.0, 1.75) {$V_3$};

      % Row labels
      \node[font=\scriptsize, gray, left=0.08cm of k1] {$K$:};
      \node[font=\scriptsize, gray, left=0.08cm of v1] {$V$:};

      % ── Arrows: each token computes its K ─────────────────────────────
      \draw[arr] (t1.south) -- (k1.north);
      \draw[arr] (t2.south) -- (k2.north);
      \draw[arr] (t3.south) -- (k3.north);

      % ── Brace: "all computed in parallel" ─────────────────────────────
      \draw[decorate, decoration={brace, amplitude=5pt, mirror}, DarkBlue]
        (0.0, 1.32) -- (5.8, 1.32);
      \node[font=\scriptsize, DarkBlue, below=0.22cm] at (2.9, 1.32)
        {Computados em \textbf{paralelo}, um único passe pelo modelo};

    \end{tikzpicture}
    \end{center}
    \vspace{0.2em}
    {\footnotesize \colorbox{CellBlue}{\,} = token \quad \colorbox{CellYellow}{\,} = computando agora}
  \end{frame}

  \begin{frame}{KV-Cache: Visualização}{Decode (Passo 1)}
    \small Gerando \texttt{"a"}: apenas $K_4,V_4$ são computados; $K_1,K_2,K_3$ vêm do cache.\\[0.3em]
    \begin{center}
    \begin{tikzpicture}[
      tok/.style={draw, rounded corners, minimum width=1.6cm, minimum height=0.52cm,
                  align=center, font=\small},
      kv/.style={draw, minimum width=1.6cm, minimum height=0.40cm,
                 align=center, font=\scriptsize},
      arr/.style={->, >=Stealth, thick},
      darr/.style={->, >=Stealth, dashed, thick, DarkBlue},
    ]
      % ── Cached prompt tokens ──────────────────────────────────────────
      \node[tok, fill=CellBlue] (t1) at (0.8, 3.5) {\texttt{O}};
      \node[tok, fill=CellBlue] (t2) at (2.9, 3.5) {\texttt{rato}};
      \node[tok, fill=CellBlue] (t3) at (5.0, 3.5) {\texttt{roeu}};
      % New token being generated
      \node[tok, fill=CellRed]  (t4) at (7.1, 3.5) {\texttt{a}};


      % ── Cached K,V (CellGreen) ────────────────────────────────────────
      \node[kv, fill=CellGreen] (k1) at (0.8, 2.45) {$K_1$};
      \node[kv, fill=CellGreen] (k2) at (2.9, 2.45) {$K_2$};
      \node[kv, fill=CellGreen] (k3) at (5.0, 2.45) {$K_3$};
      \node[kv, fill=CellGreen] (v1) at (0.8, 1.75) {$V_1$};
      \node[kv, fill=CellGreen] (v2) at (2.9, 1.75) {$V_2$};
      \node[kv, fill=CellGreen] (v3) at (5.0, 1.75) {$V_3$};

      % ── New K,V being computed (CellYellow) ───────────────────────────
      \node[kv, fill=CellYellow] (k4) at (7.1, 2.45) {$K_4$};
      \node[kv, fill=CellYellow] (v4) at (7.1, 1.75) {$V_4$};

      % Row labels
      \node[font=\scriptsize, gray, left=0.08cm of k1] {$K$:};
      \node[font=\scriptsize, gray, left=0.08cm of v1] {$V$:};

      % ── Arrow: new token computes its K ───────────────────────────────
      \draw[arr] (t4.south) -- (k4.north);

      % ── Separator: prompt vs. generated ───────────────────────────────
      \draw[dashed, gray, thin] (6.05, 1.22) -- (6.05, 3.88);
      \node[font=\tiny, gray, above] at (6.05, 3.88) {prompt $\mid$ gerado};

      % ── Attention: Q₄ reads all cached K ─────────────────────────────
      \draw[darr] (7.1, 0.92) -- (0.8, 0.92);
      \node[font=\scriptsize, DarkBlue, below=0.08cm] at (3.95, 0.92)
        {$Q_4$ consulta $K_1, K_2, K_3, K_4$ (atenção causal)};

    \end{tikzpicture}
    \end{center}
    {\footnotesize \colorbox{CellGreen}{\,} = cache (reutilizado) \quad \colorbox{CellYellow}{\,} = computando agora}
  \end{frame}

  \begin{frame}{KV-Cache: Visualização}{Decode (Passo 2)}
    \small Gerando \texttt{"rolha"}: apenas $K_5,V_5$ são computados; cache cresce para 4 entradas.\\[0.3em]
    \begin{center}
    \begin{tikzpicture}[
      tok/.style={draw, rounded corners, minimum width=1.6cm, minimum height=0.52cm,
                  align=center, font=\small},
      kv/.style={draw, minimum width=1.6cm, minimum height=0.40cm,
                 align=center, font=\scriptsize},
      arr/.style={->, >=Stealth, thick},
      darr/.style={->, >=Stealth, dashed, thick, DarkBlue},
    ]
      % ── All cached tokens ─────────────────────────────────────────────
      \node[tok, fill=CellBlue] (t1) at (0.8, 3.5) {\texttt{O}};
      \node[tok, fill=CellBlue] (t2) at (2.9, 3.5) {\texttt{rato}};
      \node[tok, fill=CellBlue] (t3) at (5.0, 3.5) {\texttt{roeu}};
      \node[tok, fill=CellBlue] (t4) at (7.1, 3.5) {\texttt{a}};
      % New token being generated
      \node[tok, fill=CellRed]  (t5) at (9.2, 3.5) {\texttt{rolha}};


      % ── Cached K,V (CellGreen) ────────────────────────────────────────
      \node[kv, fill=CellGreen] (k1) at (0.8, 2.45) {$K_1$};
      \node[kv, fill=CellGreen] (k2) at (2.9, 2.45) {$K_2$};
      \node[kv, fill=CellGreen] (k3) at (5.0, 2.45) {$K_3$};
      \node[kv, fill=CellGreen] (k4) at (7.1, 2.45) {$K_4$};
      \node[kv, fill=CellGreen] (v1) at (0.8, 1.75) {$V_1$};
      \node[kv, fill=CellGreen] (v2) at (2.9, 1.75) {$V_2$};
      \node[kv, fill=CellGreen] (v3) at (5.0, 1.75) {$V_3$};
      \node[kv, fill=CellGreen] (v4) at (7.1, 1.75) {$V_4$};

      % ── New K,V being computed (CellYellow) ───────────────────────────
      \node[kv, fill=CellYellow] (k5) at (9.2, 2.45) {$K_5$};
      \node[kv, fill=CellYellow] (v5) at (9.2, 1.75) {$V_5$};


      % Row labels
      \node[font=\scriptsize, gray, left=0.08cm of k1] {$K$:};
      \node[font=\scriptsize, gray, left=0.08cm of v1] {$V$:};

      % ── Arrow: new token computes its K ───────────────────────────────
      \draw[arr] (t5.south) -- (k5.north);

      % ── Separator: prompt vs. generated ───────────────────────────────
      \draw[dashed, gray, thin] (6.05, 1.22) -- (6.05, 3.88);
      \node[font=\tiny, gray, above] at (6.05, 3.88) {prompt $\mid$ gerado};

      % ── Attention: Q₅ reads all cached K ─────────────────────────────
      \draw[darr] (9.2, 0.92) -- (0.8, 0.92);
      \node[font=\scriptsize, DarkBlue, below=0.08cm] at (5.0, 0.92)
        {$Q_5$ consulta $K_1, \ldots, K_5$, cache cresce a cada token gerado};

    \end{tikzpicture}
    \end{center}
    {\footnotesize \colorbox{CellGreen}{\,} = cache (reutilizado) \quad \colorbox{CellYellow}{\,} = computando agora}
  \end{frame}

  %======================================================================================
  % SECTION 3: LEIS DE ESCALA
  %======================================================================================
  \section{Leis de Escala}

  \begin{frame}{O que Significa ``Escalar'' um LLM?}
    \begin{columns}[T]
      \begin{column}{0.50\textwidth}
        Três eixos de escala:
        \begin{enumerate}
          \item \textbf{Parâmetros} $N$: tamanho do modelo
          \item \textbf{Dados} $D$: tokens de treinamento
          \item \textbf{Compute} $C$: FLOPs totais de treino{\footnotesize\ (FLOP = \textit{Floating Point Operation}, unidade de custo aritmético)}
        \end{enumerate}
        \vspace{0.4em}
        Relação aproximada:
        \[
          C \approx 6 N D
        \]
        Cada token exige $\approx$6 FLOPs por parâmetro:\\
        2 no \textit{forward} $+$ 4 no \textit{backward} \cite{kaplan2020scaling}.
      \end{column}
      \begin{column}{0.46\textwidth}
        \begin{alertblock}{Pergunta central}
          Dado um orçamento de compute fixo, como distribuir entre $N$ (modelo maior) e $D$ (mais dados)?
        \end{alertblock}
        \vspace{0.4em}
        Isso é exatamente o que as \textbf{leis de escala} respondem.
      \end{column}
    \end{columns}
  \end{frame}

  \begin{frame}{Leis de Escala de Kaplan et al.\ \cite{kaplan2020scaling}}
    \begin{columns}[T]
      \begin{column}{0.52\textwidth}
        \textbf{Descoberta}: a perda de teste segue \textbf{leis de potência}:
        \[
          L(N) \propto N^{-\alpha_N}, \quad
          L(D) \propto D^{-\alpha_D}
        \]
        com $\alpha_N \approx 0.076$, $\alpha_D \approx 0.095$.\\[0.5em]
        \textbf{Recomendação}: dado compute fixo, priorize \textbf{escalar $N$} e use menos dados.
      \end{column}
      \begin{column}{0.44\textwidth}
        \begin{itemize}
          \item Válido para modelos de 768 a 1.5B parâmetros (não-embedding)
          \item Curvas de perda suaves e previsíveis
          \item Arquitetura (profundidade vs.\ largura) tem impacto menor
        \end{itemize}
        \vspace{0.3em}
        \begin{alertblock}{Limitação}
          Kaplan et al.\ não otimizavam $N$ e $D$ simultaneamente; fixavam $N$ e mediam $L$ ótima em $C$.
        \end{alertblock}
      \end{column}
    \end{columns}
  \end{frame}

  \begin{frame}{Chinchilla: Treinamento Compute-Ótimo \cite{hoffmann2022chinchilla}}
    \begin{columns}[T]
      \begin{column}{0.52\textwidth}
        \textbf{Hoffmann et al.\ (2022)}, DeepMind\\[0.4em]
        Treinaram 400+ modelos variando $N$ e $D$ com mesmo $C$.\\[0.4em]
        \begin{block}{Resultado central}
          Para compute ótimo:
          \[
            N^* \propto C^{0.5}, \quad D^* \propto C^{0.5}
          \]
          \textbf{Dobrar compute $\to$ dobrar $N$ E dobrar $D$.}
        \end{block}
      \end{column}
      \begin{column}{0.44\textwidth}
        Chinchilla (70B) $\approx$ Gopher (280B) em qualidade, com $4\times$ menos compute de inferência.\\[0.5em]
        \begin{exampleblock}{Regra prática}
          Para um modelo de $N$ parâmetros, treine com pelo menos \textbf{$20N$ tokens}.
        \end{exampleblock}
        \vspace{0.3em}
        Exemplo: modelo de 7B $\to$ $\ge$140B tokens.
      \end{column}
    \end{columns}
  \end{frame}

  \begin{frame}{Relação Ótima Tokens/Parâmetros}
    \begin{columns}[T]
      \begin{column}{0.52\textwidth}
        \resizebox{\linewidth}{!}{%
        \begin{tikzpicture}
          \begin{axis}[
            xlabel={Parâmetros (B)},
            ylabel={Tokens ótimos (B)},
            xmin=1, xmax=140,
            ymin=0, ymax=3000,
            grid=both,
            grid style={line width=0.3pt, draw=gray!30},
            width=7.5cm, height=5.5cm,
            thick
          ]
            \addplot[DarkBlue, thick, domain=1:140, samples=50]
              {20*x};
            \addplot[only marks, mark=*, mark size=2.5pt, DarkRed]
              coordinates {(7,2000) (13,2000) (70,2000)};
            \node[font=\tiny, left] at (axis cs:7,2000) {7B};
            \node[font=\tiny, above right] at (axis cs:13,2000) {13B};
            \node[font=\tiny, above right] at (axis cs:70,2000) {70B};
            \legend{$20N$ (Chinchilla), Llama 2 (2T)}
          \end{axis}
        \end{tikzpicture}}
      \end{column}
      \begin{column}{0.44\textwidth}
        \begin{itemize}
          \item Linha azul: regra dos $20N$ tokens (compute-ótimo, Chinchilla)
          \item \textbf{Llama 2} usou 2T tokens em todos os tamanhos, bem acima do $20N$ (sobretreino proposital)
          \item Tendência atual: ainda mais dados. Llama~3 8B usou $15$T, $\sim\!93\times$ o $20N$ ($=160$B)
        \end{itemize}
        \vspace{0.3em}
        \begin{alertblock}{Compute-ótimo $\ne$ inferência-ótimo}
          Modelos menores treinados por mais tempo custam menos na inferência, preferíveis quando há muitas chamadas.
        \end{alertblock}
      \end{column}
    \end{columns}
  \end{frame}

  \begin{frame}{Implicações Práticas das Leis de Escala}
    \begin{columns}[T]
      \begin{column}{0.50\textwidth}
        \begin{itemize}
          \item \textbf{Llama 1} (Meta, 2023): 7B--65B treinados com 1T--1.4T tokens
          \item \textbf{Llama 2} (Meta, 2023): até 2T tokens, contexto 4096
          \item \textbf{Llama 3} (Meta, 2024): 8B e 70B com 15T tokens
          \item \textbf{Mistral 7B} (2023): supera Llama 2 13B com 7B parâmetros
        \end{itemize}
      \end{column}
      \begin{column}{0.46\textwidth}
        \begin{exampleblock}{Lição aprendida}
          Treinar modelos \textbf{menores por mais tempo} é mais custo-efetivo para implantação do que modelos gigantes subtreinados.
        \end{exampleblock}
        \vspace{0.3em}
        A Chinchilla reorientou o campo: de corrida por parâmetros (GPT-3, Gopher, Megatron) para equilíbrio $N \times D$.
      \end{column}
    \end{columns}
  \end{frame}

  \begin{frame}{Modelos Recentes: Parâmetros, Dados e Contexto}
    \resizebox{\linewidth}{!}{%
    \begin{tabular}{lrrrrr}
      \toprule
      \textbf{Modelo} & \textbf{Params} & \textbf{Tokens treino} & \textbf{Contexto} & \textbf{Norm} & \textbf{Atenção} \\
      \midrule
      GPT-2 XL       & 1.5B   & $\sim$40B      & 1\,024  & Pre-LN  & MHA \\
      GPT-3          & 175B   & 300B           & 2\,048  & Pre-LN  & MHA \\
      Llama 2 7B     & 7B     & 2T             & 4\,096  & RMSNorm & GQA \\
      Llama 2 70B    & 70B    & 2T             & 4\,096  & RMSNorm & GQA \\
      Llama 3 8B     & 8B     & 15T            & 8\,192  & RMSNorm & GQA \\
      Mistral 7B     & 7B     & $\sim$1T       & 32\,768 & RMSNorm & GQA+SWA \\
      Mixtral 8×7B   & 46.7B  & $\sim$1T       & 32\,768 & RMSNorm & GQA+MoE \\
      \bottomrule
    \end{tabular}}
    \vspace{0.3em}
    \begin{center}
      \footnotesize SWA = Sliding Window Attention; MoE = Mistura de Especialistas
    \end{center}
  \end{frame}

  %======================================================================================
  % SECTION 5: APROXIMAÇÕES DE ATENÇÃO
  %======================================================================================
  \section{Aproximações de Atenção}

  \begin{frame}{O Gargalo Quadrático da Atenção}
    \begin{columns}[T]
      \begin{column}{0.50\textwidth}
        \textbf{Custo da atenção padrão} \cite{vaswani2017attention}:
        \begin{itemize}
          \item Memória: $\mathcal{O}(n^2)$, matriz de scores
          \item Compute: $\mathcal{O}(n^2 d)$, produto $QK^\top$
        \end{itemize}
        \vspace{0.3em}
        Com $n=32\,768$ tokens e $H=32$ cabeças:\\
        Matriz de atenção: $32768^2 \times 32 \approx 34$ bilhões de valores $\approx 69$ GB (float16)
      \end{column}
      \begin{column}{0.46\textwidth}
        \resizebox{\linewidth}{!}{%
        \begin{tikzpicture}
          \begin{axis}[
            xlabel={Comprimento $n$},
            ylabel={Custo relativo},
            xmin=512, xmax=32768,
            xmode=log, log basis x=2,
            ymode=log,
            ymin=1, ymax=5000,
            xtick={512,2048,8192,32768},
            xticklabels={512,2K,8K,32K},
            grid=both,
            grid style={line width=0.3pt,draw=gray!30},
            width=7cm, height=4.5cm,
            thick
          ]
            \addplot[DarkRed, thick, domain=512:32768, samples=40]
              {(x/512)^2};
            \addplot[DarkBlue, thick, domain=512:32768, samples=10]
              {(x/512)};
            \legend{$\mathcal{O}(n^2)$, $\mathcal{O}(n)$}
          \end{axis}
        \end{tikzpicture}}
      \end{column}
    \end{columns}
    \vspace{0.3em}
    \textcolor{DarkRed}{\textbf{Motivação:}} reduzir custo ao escalar o contexto \cite{stanfordcme295}.
  \end{frame}

  \begin{frame}{Multi-Head Attention (MHA)}{Recapitulação}
    \begin{columns}[T]
      \begin{column}{0.52\textwidth}
        Com $H$ cabeças, $d_k = d_\text{model}/H$:
        \begin{itemize}
          \item $H$ projeções de Q, K, V independentes
          \item $H$ pares de matrizes KV em cache por camada
          \item Custo de KV-cache por token: $2 \cdot H \cdot d_k \cdot L$
        \end{itemize}
        \vspace{0.3em}
        Para GPT-3 ($H=96$, $d_k=128$, $L=96$):
        \[
          \text{Cache} = 2 \times 96 \times 128 \times 96 \approx 2.4\text{M floats/token}
        \]
      \end{column}
      \begin{column}{0.44\textwidth}
        \resizebox{\linewidth}{!}{%
        \begin{tikzpicture}[
          head/.style={draw, fill=CellBlue, minimum width=1.1cm, minimum height=1.4cm,
                       rounded corners, font=\small, align=center},
          node distance=0.4cm
        ]
          \foreach \i/\l in {1/1,2/2,3/3,4/H} {
            \node[head] (h\i) at (\i*1.3, 0) {Q\textsubscript{\l}\\K\textsubscript{\l}\\V\textsubscript{\l}};
          }
          \node[font=\small] at (4*1.3+0.9, 0) {\ldots};
          \node[font=\small] at (2.6, -1.1) {KV-cache: $H$ pares por camada};
        \end{tikzpicture}}
        \vspace{0.3em}
        Gargalo em inferência: KV-cache cresce com $n$, $H$, $L$.
      \end{column}
    \end{columns}
  \end{frame}

  \begin{frame}{Multi-Query Attention (MQA) \cite{shazeer2019mqa}}
    \begin{columns}[T]
      \begin{column}{0.52\textwidth}
        \textbf{Shazeer (2019)}: reduza K e V para \textbf{uma única cabeça} compartilhada; mantenha $H$ cabeças de Q.
        \[
          \text{head}_h = \text{Attn}(QW_h^Q,\; K W^K,\; V W^V)
        \]
        \textbf{Redução de cache:} $H\times \to 1\times$ nos tensores K, V.\\[0.3em]
        Cache por token: $2 \cdot d_k \cdot L$ (independente de $H$).
      \end{column}
      \begin{column}{0.44\textwidth}
        \resizebox{\linewidth}{!}{%
        \begin{tikzpicture}[
          qhead/.style={draw, fill=CellBlue, minimum width=0.8cm, minimum height=0.85cm,
                        rounded corners, font=\small, align=center},
          kvhead/.style={draw, fill=CellGreen, minimum width=1.0cm, minimum height=0.85cm,
                         rounded corners, font=\small, align=center},
          arrow/.style={->, >=Stealth, thick}
        ]
          \foreach \i in {1,2,3,4} {
            \node[qhead] (q\i) at (\i*0.95, 1.5) {Q\textsubscript{\i}};
          }
          \node[kvhead] (kv) at (2.375, 0) {K\,V};
          \foreach \i in {1,2,3,4} {
            \draw[arrow] (kv) -- (q\i);
          }
          \node[right=0.2cm of q4, font=\footnotesize] {$H$ cabeças Q};
          \node[below=0.12cm of kv, font=\footnotesize] {1 par K,V};
        \end{tikzpicture}}
        \begin{itemize}
          \item KV-cache $H\times$ menor
          \item Inferência muito mais rápida
          \item Leve queda de qualidade
        \end{itemize}
      \end{column}
    \end{columns}
  \end{frame}

  \begin{frame}{Grouped Query Attention (GQA) \cite{ainslie2023gqa}}
    \begin{columns}[T]
      \begin{column}{0.52\textwidth}
        \textbf{Ainslie et al.\ (2023)}: compromisso entre MHA e MQA.\\[0.4em]
        Divida as $H$ cabeças de Q em $G$ grupos. Cada grupo compartilha um par K,V.
        \[
          G = 1 \;\Rightarrow\; \text{MQA}
          \qquad
          G = H \;\Rightarrow\; \text{MHA}
        \]
        Llama 2/3, Mistral, Mixtral usam GQA.
      \end{column}
      \begin{column}{0.44\textwidth}
        \centering
        \begin{tikzpicture}[
          qhead/.style={draw, fill=CellBlue, minimum width=0.75cm, minimum height=0.8cm,
                        rounded corners, font=\small, align=center},
          kvhead/.style={draw, fill=CellGreen, minimum width=0.9cm, minimum height=0.8cm,
                         rounded corners, font=\small, align=center},
          arrow/.style={->, >=Stealth, thick}
        ]
          \node[qhead] (q1) at (0.5,1.6) {Q\textsubscript{1}};
          \node[qhead] (q2) at (1.3,1.6) {Q\textsubscript{2}};
          \node[kvhead] (kva) at (0.9,0) {K\textsubscript{a}V\textsubscript{a}};
          \draw[arrow] (kva)--(q1); \draw[arrow] (kva)--(q2);

          \node[qhead] (q3) at (2.8,1.6) {Q\textsubscript{3}};
          \node[qhead] (q4) at (3.6,1.6) {Q\textsubscript{4}};
          \node[kvhead] (kvb) at (3.2,0) {K\textsubscript{b}V\textsubscript{b}};
          \draw[arrow] (kvb)--(q3); \draw[arrow] (kvb)--(q4);

          \node[font=\footnotesize, below=0.08cm of kva] {grupo 1};
          \node[font=\footnotesize, below=0.08cm of kvb] {grupo 2};
        \end{tikzpicture}
        \begin{itemize}
          \item $G=8$ típico (Llama 2 70B: $H=64$, $G=8$)
          \item Cache reduzido $H/G\times$
          \item Qualidade quase idêntica ao MHA
        \end{itemize}
      \end{column}
    \end{columns}
  \end{frame}

  \begin{frame}{Comparação: MHA, MQA e GQA}
    \begin{center}
    \resizebox{0.92\linewidth}{!}{%
    \begin{tabular}{lcccl}
      \toprule
      \textbf{Variante} & \textbf{Pares K,V} & \textbf{Cache KV} & \textbf{Qualidade} & \textbf{Exemplos} \\
      \midrule
      MHA & $H$            & $H \cdot d_k$ & Referência         & GPT-2, GPT-3, BERT \\
      MQA & $1$            & $d_k$         & $\downarrow$ levemente & PaLM, Falcon \\
      GQA & $G$ ($1<G<H$) & $G \cdot d_k$ & $\approx$ MHA     & Llama 2/3, Mistral \\
      \bottomrule
    \end{tabular}}
    \end{center}
    \vspace{0.4em}
    \begin{columns}[T]
      \begin{column}{0.50\textwidth}
        \begin{exampleblock}{GQA na prática}
          Llama 2 70B: $H=64$ cabeças Q, $G=8$ grupos KV.\\
          Cache $8\times$ menor que MHA com qualidade semelhante.
        \end{exampleblock}
      \end{column}
      \begin{column}{0.46\textwidth}
        \begin{alertblock}{Por que importa?}
          Redução do KV-cache é fundamental para servir muitos usuários simultaneamente, o cache de todos os requests cabe na GPU.
        \end{alertblock}
      \end{column}
    \end{columns}
  \end{frame}

  \begin{frame}{Sliding Window Attention \cite{jiang2023mistral}}
    \begin{columns}[T]
      \begin{column}{0.52\textwidth}
        \textbf{Motivação}: em sequências longas, a maioria das informações úteis está \textbf{próxima} do token atual.\\[0.4em]
        Cada posição $i$ atende apenas às $W$ posições anteriores:
        \[
          \text{Attn}(i) \leftarrow \text{tokens em } [i-W,\; i]
        \]
        Custo: $\mathcal{O}(n \cdot W)$ em vez de $\mathcal{O}(n^2)$.\\[0.3em]
        Campo receptivo em $L$ camadas: $L \cdot W$.\\
        {\footnotesize Informação distante propaga-se por composição: camada $l$ agrega o que camada $l{-}1$ já resumiu da janela anterior.}
      \end{column}
      \begin{column}{0.44\textwidth}
        \begin{block}{Configuração Mistral 7B}
          \begin{itemize}
            \item Janela local: $W = 4096$
            \item 32 camadas
            \item Campo receptivo: $32 \times 4096 = 131\text{K}$
            \item Contexto declarado: 32K tokens
          \end{itemize}
        \end{block}
        \vspace{0.3em}
        \textcolor{DarkBlue}{\textbf{Resultado}}: Mistral 7B supera Llama 2 13B em benchmarks com $2\times$ menos parâmetros.
      \end{column}
    \end{columns}
  \end{frame}

  \begin{frame}{Resumo: Variantes de Atenção e Trade-offs}
    \begin{center}
    \resizebox{\linewidth}{!}{%
    \begin{tabular}{lccccc}
      \toprule
      \textbf{Variante} & \textbf{Compute} & \textbf{Mem.\ KV} & \textbf{Contexto longo} & \textbf{Qualidade} & \textbf{Exemplos} \\
      \midrule
      MHA (padrão) & $\mathcal{O}(n^2 d)$ & $H \cdot d_k$ & Limitado    & Referência           & GPT-2/3 \\
      MQA          & $\mathcal{O}(n^2 d)$ & $d_k$         & Moderado    & $\downarrow$ levemente & PaLM, Falcon \\
      GQA          & $\mathcal{O}(n^2 d)$ & $G \cdot d_k$ & Moderado    & $\approx$ MHA        & Llama 2/3, Mistral \\
      SWA          & $\mathcal{O}(n W d)$ & $W \cdot d_k$ & Alto        & Boa (local)          & Mistral \\
      \bottomrule
    \end{tabular}}
    \end{center}
    \vspace{0.2em}
    \begin{center}
      \footnotesize Baseado em \cite{stanfordcme295}. SWA = Sliding Window Attention.
    \end{center}
  \end{frame}

  %======================================================================================
  % SECTION 6: MISTURA DE ESPECIALISTAS
  %======================================================================================
  \section{Mistura de Especialistas (MoE)}

  \begin{frame}{Motivação: Mais Parâmetros sem Mais Compute}
    \begin{columns}[T]
      \begin{column}{0.52\textwidth}
        \textbf{Dilema}:
        \begin{itemize}
          \item Modelos maiores $\to$ melhor qualidade
          \item Modelos maiores $\to$ mais compute por token
        \end{itemize}
        \vspace{0.4em}
        \textbf{Ideia da Mistura de Especialistas (MoE)}:\\
        Substitua a FFN densa por $E$ redes menores (\textit{especialistas}).\\
        Cada token ativa apenas $k$ especialistas ($k \ll E$).
      \end{column}
      \begin{column}{0.44\textwidth}
        \begin{exampleblock}{Exemplo: Mixtral 8×7B}
          8 especialistas FFN\\
          Apenas 2 ativados por token\\
          $\to$ 46.7B parâmetros totais\\
          $\to$ compute $\approx$ 12.9B por token
        \end{exampleblock}
        \vspace{0.3em}
        Origem: \textbf{Sparsely-Gated MoE} \cite{shazeer2017moe} (2017), revivido em LLMs a partir de 2022.
      \end{column}
    \end{columns}
  \end{frame}

  \begin{frame}{Arquitetura MoE: Router + Especialistas}
    \begin{columns}[T]
      \begin{column}{0.46\textwidth}
        \centering
        \begin{tikzpicture}[
          router/.style={draw, fill=CellYellow, rounded corners,
                         minimum width=3cm, minimum height=0.7cm, font=\small, align=center},
          expert/.style={draw, fill=CellBlue, rounded corners,
                         minimum width=1.6cm, minimum height=0.65cm, font=\small, align=center},
          sum/.style={draw, circle, fill=white, minimum size=0.45cm, font=\small},
          arrow/.style={->, >=Stealth, thick},
          node distance=0.5cm
        ]
          \node[router] (r) {Router $G(x)$};
          \node[expert, above left=0.9cm and 0.2cm of r] (e1) {FFN$_1$};
          \node[expert, above=0.9cm of r]               (e2) {FFN$_2$};
          \node[expert, above right=0.9cm and 0.2cm of r] (e3) {FFN$_3$};
          \node[sum, above=2.7cm of r] (s) {$\sum$};
          \node[above=0.3cm of s, font=\small] (out) {saída};
          \node[below=0.5cm of r, font=\small] (inp) {entrada $x$};

          \draw[arrow] (inp) -- (r);
          \draw[arrow] (r) -- (e1);
          \draw[arrow] (r) -- (e2);
          \draw[arrow] (r) -- (e3);
          \draw[arrow] (e1) |- (s);
          \draw[arrow] (e2) -- (s);
          \draw[arrow] (e3) |- (s);
          \draw[arrow] (s) -- (out);
        \end{tikzpicture}
      \end{column}
      \begin{column}{0.50\textwidth}
        \textbf{Em cada bloco MoE}:
        \begin{enumerate}
          \item Router calcula scores para cada especialista
          \item Seleciona top-$k$ especialistas
          \item Processa $x$ nos $k$ especialistas
          \item Combina saídas ponderadas pelos scores
        \end{enumerate}
        \vspace{0.4em}
        \begin{block}{Nota importante}
          As camadas de \textbf{atenção} permanecem densas.\\
          Somente a FFN vira MoE.
        \end{block}
      \end{column}
    \end{columns}
  \end{frame}

  \begin{frame}{Top-$k$ Routing: Mecanismo de Gating}
    \begin{columns}[T]
      \begin{column}{0.52\textwidth}
        Dado token $x \in \mathbb{R}^d$ e pesos $W_g \in \mathbb{R}^{d \times E}$:
        \[
          g(x) = \text{softmax}\!\left(\text{TopK}\!\left(x W_g,\; k\right)\right)
        \]
        onde $\text{TopK}$ mantém os $k$ maiores scores e zera o resto.\\[0.4em]
        Saída:
        \[
          y = \sum_{i \in \text{Top-}k} g_i(x)\; \text{FFN}_i(x)
        \]
      \end{column}
      \begin{column}{0.44\textwidth}
      \scriptsize
        \begin{itemize}
          \item Mixtral usa $k=2$, $E=8$
          \item $\sim$25\% dos parâmetros ativos por token
          \item Backward: apenas pelos $k$ especialistas ativos
        \end{itemize}
        \vspace{0.3em}
        \begin{alertblock}{Expert collapse}
          \textbf{Causa}: o router aprende a enviar tokens ao especialista que já é o melhor, que melhora ainda mais, atraindo mais tokens, em ciclo de reforço positivo.\\[0.15em]
          \textbf{Resultado}: 1--2 especialistas processam quase tudo; os demais recebem gradiente zero e nunca aprendem.\\[0.15em]
          \textbf{Solução}: \textit{Loss} auxiliar de balanceamento (próximo slide).
        \end{alertblock}
      \end{column}
    \end{columns}
  \end{frame}

  \begin{frame}{Problemas de MoE: Balanceamento de Carga}
    \begin{columns}[T]
      \begin{column}{0.50\textwidth}
        \textbf{Perda auxiliar de balanceamento}:
        \[
          \mathcal{L}_\text{aux} = \alpha \cdot E \cdot \sum_{i=1}^E f_i \cdot P_i
        \]
        onde $f_i$ = fração de tokens roteados ao especialista $i$,\\
        $P_i$ = score médio do router para o especialista $i$.\\[0.3em]
        Minimizar $\mathcal{L}_\text{aux}$ encoraja distribuição uniforme.
      \end{column}
      \begin{column}{0.46\textwidth}
        Outros desafios:
        \begin{itemize}
          \item \textbf{Instabilidade}: gradientes esparsas podem causar divergência
          \item \textbf{Comunicação}: especialistas distribuídos em GPUs diferentes (\textit{model parallelism})
          \item \textbf{Capacidade máxima}: tokens descartados se um especialista ficar sobrecarregado
        \end{itemize}
        \vspace{0.2em}
        \begin{exampleblock}{Mixtral}
          Usa $\alpha=0.01$, roteamento token a token.
        \end{exampleblock}
      \end{column}
    \end{columns}
  \end{frame}

  \begin{frame}{Mixtral 8×7B \cite{jiang2024mixtral}}
    \begin{columns}[T]
      \begin{column}{0.52\textwidth}
        \textbf{Jiang et al.\ (2024)}, Mistral AI\\[0.4em]
        Arquitetura:
        \begin{itemize}
          \item 32 camadas transformer
          \item Cada FFN $\to$ 8 especialistas de 14336 neurônios
          \item Top-2 routing por token
          \item GQA ($G=8$), RMSNorm, SWA
          \item Contexto: 32\,768 tokens
        \end{itemize}
      \end{column}
      \begin{column}{0.44\textwidth}
        \resizebox{\linewidth}{!}{%
        \begin{tabular}{lrr}
          \toprule
          \textbf{Modelo} & \textbf{Params ativos} & \textbf{Params totais} \\
          \midrule
          Llama 2 13B   & 13B  & 13B   \\
          Llama 2 70B   & 70B  & 70B   \\
          Mistral 7B    & 7B   & 7B    \\
          Mixtral 8×7B  & 12.9B & 46.7B \\
          \bottomrule
        \end{tabular}}
        \vspace{0.3em}
        \begin{exampleblock}{Resultado}
          Mixtral 8×7B supera Llama 2 70B em benchmarks com $5.4\times$ menos parâmetros ativos.
        \end{exampleblock}
      \end{column}
    \end{columns}
  \end{frame}

  \begin{frame}{MoE Sparse vs.\ Modelos Densos}
    \begin{center}
    \resizebox{0.90\linewidth}{!}{%
    \begin{tabular}{lcc}
      \toprule
      \textbf{Característica} & \textbf{Denso (Llama 2 70B)} & \textbf{Sparse MoE (Mixtral 8×7B)} \\
      \midrule
      Parâmetros totais       & 70B   & 46.7B \\
      Parâmetros ativos       & 70B   & 12.9B \\
      Compute por token       & Alto  & $\sim5\times$ menor \\
      Memória total (fp16)    & $\sim$140 GB & $\sim$93 GB \\
      Throughput (tokens/s)   & Menor & Maior \\
      Qualidade (MMLU)        & 68.9\% & 70.6\% \\
      \bottomrule
    \end{tabular}}
    \end{center}
    \vspace{0.4em}
    \begin{center}
      \textcolor{DarkBlue}{\textbf{MoE: mais capacidade com menor custo de inferência}},\\
      ao custo de maior memória total e complexidade de treino.
    \end{center}
  \end{frame}

  %======================================================================================
  % SECTION 7: DECODIFICAÇÃO
  %======================================================================================
  \section{Estratégias de Decodificação}

  \begin{frame}{Estratégias de Decodificação: Greedy, Beam, Sampling}
    \begin{columns}[T]
      \begin{column}{0.50\textwidth}
        Após o LM produzir logits $z \in \mathbb{R}^{|V|}$:
        \begin{block}{Greedy}
          $x_t = \arg\max_v P(v \mid x_{<t})$\\
          Rápido, determinístico, pode ser subótimo.
        \end{block}
        \begin{block}{Beam Search}
          Mantém $B$ hipóteses; escolhe a de maior probabilidade conjunta.
          $\mathcal{O}(B)$ mais lento; ótimo para tradução.
        \end{block}
      \end{column}
      \begin{column}{0.46\textwidth}
        \begin{block}{Amostragem}
          $x_t \sim P(\cdot \mid x_{<t})$\\
          Estocástico; necessário para criatividade.
        \end{block}
        \vspace{0.3em}
        Parâmetros de controle:
        \begin{itemize}
          \item \textbf{Temperature} $T$: $P_T \propto \exp(z_v / T)$
          \item \textbf{Top-$k$}: restringe ao top-$k$ tokens
          \item \textbf{Top-$p$} (nucleus): restringe ao menor conjunto com prob.\ acumulada $\ge p$
        \end{itemize}
      \end{column}
    \end{columns}
  \end{frame}

  \begin{frame}{Temperature, Top-$k$ e Top-$p$}
    \begin{columns}[T]
      \begin{column}{0.52\textwidth}
        \textbf{Temperature}:
        \[
          P_T(x) = \frac{\exp(z_x / T)}{\sum_v \exp(z_v / T)}
        \]
        \begin{itemize}
          \item $T \to 0$: greedy
          \item $T = 1$: distribuição original
          \item $T > 1$: distribuição mais uniforme
        \end{itemize}
        \vspace{0.3em}
        \textbf{Top-$p$}: selecione o menor conjunto de tokens cuja probabilidade acumulada $\ge p$ (ex: $p=0.9$), então amostre dentro desse conjunto.
      \end{column}
      \begin{column}{0.44\textwidth}
        \resizebox{\linewidth}{!}{%
        \begin{tikzpicture}
          \begin{axis}[
            ybar,
            bar width=0.38cm,
            xlabel={Token},
            ylabel={Probabilidade},
            symbolic x coords={$v_1$,$v_2$,$v_3$,$v_4$,$v_5$,$v_6$,$v_7$,$v_8$},
            xtick=data,
            ymin=0, ymax=0.55,
            width=7cm, height=4.5cm,
            enlarge x limits=0.08
          ]
            \addplot[fill=CellBlue] coordinates
              {($v_1$,0.45)($v_2$,0.25)($v_3$,0.15)($v_4$,0.07)($v_5$,0.04)($v_6$,0.02)($v_7$,0.01)($v_8$,0.01)};
            \draw[DarkRed, dashed, thick]
              (axis cs:{$v_4$}, 0) -- (axis cs:{$v_4$}, 0.55)
              node[above, font=\tiny, DarkRed] {top-$p{=}0.85$};
          \end{axis}
        \end{tikzpicture}}
      \end{column}
    \end{columns}
  \end{frame}

  \begin{frame}{Decodificação Guiada (\textit{Constrained Decoding}) \cite{willard2023guided}}
    \begin{columns}[T]
      \begin{column}{0.52\textwidth}
        \textbf{Problema}: LLMs podem gerar texto que não segue formato específico (JSON, SQL, código).\\[0.4em]
        \textbf{Solução}: restringir a distribuição de saída a tokens válidos segundo uma \textbf{gramática} ou \textbf{esquema}.\\[0.4em]
        \textbf{Willard \& Louf (2023)}: pré-compila um DFA para cada gramática, o mascaramento de tokens inválidos é $\mathcal{O}(1)$ por token.
      \end{column}
      \begin{column}{0.44\textwidth}
        \begin{block}{Exemplos de uso}
          \begin{itemize}
            \item Geração de JSON estruturado
            \item SQL sem erros de sintaxe
            \item Extração em formato fixo
            \item Chamadas de API com parâmetros corretos
          \end{itemize}
        \end{block}
        \vspace{0.2em}
        Implementações: \texttt{outlines}, \texttt{guidance}, \texttt{llama.cpp} grammar mode, vLLM structured outputs.
      \end{column}
    \end{columns}
  \end{frame}

  \begin{frame}{Decodificação Especulativa: Ideia Geral \cite{leviathan2023speculative}}
    \begin{columns}[T]
      \begin{column}{0.52\textwidth}
        \textbf{Gargalo}: LLMs grandes são lentos, cada token exige um passe completo.\\[0.4em]
        \textbf{Observação}: a maioria dos tokens é fácil de prever; tokens difíceis são exceção.\\[0.4em]
        \textbf{Algoritmo}:
        \begin{enumerate}
          \item \textbf{Draft model} pequeno gera $\gamma$ tokens
          \item \textbf{Target model} verifica \textit{em paralelo}
          \item No primeiro rejeitado: re-amostra o token e \textbf{descarta o restante} do rascunho
          \item Todos aceitos: amostra um token extra do passe do alvo, ganhando $\gamma + 1$ tokens
        \end{enumerate}
      \end{column}
      \begin{column}{0.44\textwidth}
        \begin{block}{Propriedade chave}
          Saída \textbf{idêntica em distribuição} ao modelo alvo, sem perda de qualidade por construção.
        \end{block}
        \vspace{0.3em}
        \begin{exampleblock}{Speedup típico}
          $2$--$3\times$ com draft $10\times$ menor.\\
          Ex: Llama 3 70B + Llama 3 8B como draft.
        \end{exampleblock}
      \end{column}
    \end{columns}
  \end{frame}

  \begin{frame}{Algoritmo Especulativo: Critério de Aceitação}
    \begin{columns}[T]
      \begin{column}{0.52\textwidth}
        Draft model $M_q$, target model $M_p$:\\[0.4em]
        Para cada token $\tilde{x}_t$ proposto pelo draft:
        \[
          \text{Aceitar com prob.} \min\!\left(1,\; \frac{p(\tilde{x}_t \mid x_{<t})}{q(\tilde{x}_t \mid x_{<t})}\right)
        \]
        Se rejeitado: re-amostre de
        \[
          p'(x) \propto \max\!\bigl(0,\, p(x)-q(x)\bigr)
        \]
        Verifique todos os $\gamma$ tokens \textbf{em um único passe} do modelo alvo.
      \end{column}
      \begin{column}{0.44\textwidth}
        \begin{itemize}
          \item $p \ge q$: token \textbf{sempre aceito}
          \item Draft e target concordam $\to$ aceleração máxima
          \item Tokens difíceis: target corrige
          \item Rejeição na posição $t$: as posições $t{+}1, \ldots, \gamma$ do rascunho são \textbf{descartadas}
        \end{itemize}
        \vspace{0.3em}
        \begin{block}{Complexidade}
          $\gamma+1$ passos de $M_q$ + 1 passe de $M_p$\\
          $\to$ até $\gamma+1$ tokens gerados\\
          ($\gamma$ aceitos + 1 extra amostrado de $M_p$).\\
          Sem especulação: $\gamma+1$ passes de $M_p$.
        \end{block}
      \end{column}
    \end{columns}
  \end{frame}

  \begin{frame}{Decodificação Especulativa: Exemplo Passo a Passo}
    \textbf{Contexto:} completar \enquote{O gato subiu na \_\_\_} com $\gamma=3$.\\[0.3em]
    Draft $M_q$ propõe: \texttt{árvore}, \texttt{e}, \texttt{dormiu}, verificados em \textbf{um único passe} de $M_p$.\\[0.4em]
    \begin{columns}[T]
      \begin{column}{0.58\textwidth}
      \resizebox{\linewidth}{!}{
        \begin{tabular}{lcccc}
          \toprule
          Token & $q$ & $p$ & $\min\!\left(1,\tfrac{p}{q}\right)$ & Decisão \\
          \midrule
          \texttt{árvore} & 0.6 & 0.9 & \textbf{1,00} & sempre aceito \\
          \texttt{e}      & 0.7 & 0.4 & 0.57 & aceito c/ prob.\ 57\% \\
          \texttt{dormiu} & 0.8 & 0.2 & 0.25 & aceito c/ prob.\ 25\% \\
          \bottomrule
        \end{tabular}}
      \end{column}
      \begin{column}{0.38\textwidth}
        \begin{exampleblock}{Cenário A (todos aceitos)}
          3 aceitos + 1 extra de $M_p$:\\
          4 tokens em 1 passe.\\
          Sem especulação: 4 passes.
        \end{exampleblock}
        \vspace{0.3em}
        \begin{alertblock}{Cenário B (\texttt{e} rejeitado)}
          Re-amostra \texttt{e} de $p'\!\propto\!\max(0,p{-}q)$ e descarta \texttt{dormiu};\\
          \texttt{árvore} + token corrigido, ainda 1 passe.
        \end{alertblock}
      \end{column}
    \end{columns}
  \end{frame}

  %─── Slide: Rejeição ─────────────────────────────────────────────────────────
  \begin{frame}{Rejeição: Distribuições $q$ e $p$ para \texttt{dormiu}}
    \textbf{Contexto:} ``O gato subiu na árvore e \_\_\_''
    (\texttt{árvore} e \texttt{e} já aceitos).\\[0.3em]
    Draft $M_q$ propõe \texttt{dormiu} com $q{=}0.8$; modelo alvo tem $p{=}0.2$.\\[0.2em]
    \begin{columns}[T]
      \begin{column}{0.52\textwidth}
        \resizebox{\linewidth}{!}{%
        \begin{tikzpicture}
          \begin{axis}[
            ybar,
            bar width=0.30cm,
            ylabel={Probabilidade},
            xlabel={Token candidato},
            symbolic x coords={dormiu,pulou,caiu,desceu},
            xtick=data,
            x tick label style={font=\small\ttfamily},
            ymin=0, ymax=0.95,
            ytick={0,0.2,0.4,0.6,0.8},
            width=7.5cm, height=5.0cm,
            enlarge x limits=0.18,
            legend style={at={(0.98,0.97)}, anchor=north east, font=\small},
            legend cell align=left,
          ]
            \addplot[fill=CellBlue, draw=DarkBlue] coordinates {
              (dormiu,0.80)(pulou,0.10)(caiu,0.05)(desceu,0.05)
            };
            \addlegendentry{$q$ (draft $M_q$)}
            \addplot[fill=CellYellow, draw=DarkRed] coordinates {
              (dormiu,0.20)(pulou,0.40)(caiu,0.25)(desceu,0.15)
            };
            \addlegendentry{$p$ (alvo $M_p$)}
            \node[above, font=\scriptsize\bfseries]
              at (axis cs:dormiu, 0.85) {\textcolor{DarkRed}{REJEITADO}};
            \draw[DarkRed, very thick, ->]
              (axis cs:dormiu, 0.83) -- (axis cs:dormiu, 0.78);
          \end{axis}
        \end{tikzpicture}}
      \end{column}
      \begin{column}{0.44\textwidth}
        \scriptsize
        \begin{alertblock}{Por que \texttt{dormiu} é rejeitado?}
          $p(\texttt{dormiu}) = 0.2 < q(\texttt{dormiu}) = 0.8$\\[0.3em]
          Aceitar com probabilidade:
          \[
            \min\!\left(1,\;\frac{p}{q}\right) = \frac{0.2}{0.8} = 0.25
          \]
          Neste cenário, \textbf{\texttt{dormiu} é rejeitado.}
        \end{alertblock}
        \vspace{0.4em}
        \begin{block}{Observação}
          $M_p$ prefere \texttt{pulou} ($p{=}0.4$), o draft superestimou a certeza em \texttt{dormiu}.
        \end{block}
      \end{column}
    \end{columns}
  \end{frame}

  %─── Slide: Re-amostragem ────────────────────────────────────────────────────
  \begin{frame}{Re-amostragem: $p'(x) \propto \max\!\bigl(0,\, p(x) - q(x)\bigr)$}
    \begin{columns}[T]
      \begin{column}{0.50\textwidth}
        \resizebox{\linewidth}{!}{%
        \begin{tabular}{lcccc}
          \toprule
          Token & $p$ & $q$ & $p{-}q$ & $p'$ (norm.) \\
          \midrule
          \texttt{dormiu} & 0.20 & 0.80 & \textcolor{DarkRed}{$-$0.60} & \cellcolor{CellRed}0.00 \\
          \texttt{pulou}  & 0.40 & 0.10 & $+$0.30 & \cellcolor{CellGreen}\textbf{0.50} \\
          \texttt{caiu}   & 0.25 & 0.05 & $+$0.20 & \cellcolor{CellGreen}0.33 \\
          \texttt{desceu} & 0.15 & 0.05 & $+$0.10 & \cellcolor{CellGreen}0.17 \\
          \bottomrule
        \end{tabular}}
        \vspace{0.4em}
        \begin{exampleblock}{Token selecionado de $p'$}
          \textbf{\texttt{pulou}} ($p'{=}0.50$)\\
          ``O gato subiu na árvore e \textbf{\texttt{pulou}}\ldots''
        \end{exampleblock}
      \end{column}
      \begin{column}{0.46\textwidth}
        \resizebox{.8\linewidth}{!}{%
        \begin{tikzpicture}
          \begin{axis}[
            ybar,
            bar width=0.45cm,
            ylabel={$p'(x)$ normalizado},
            xlabel={Token candidato},
            symbolic x coords={dormiu,pulou,caiu,desceu},
            xtick=data,
            x tick label style={font=\small\ttfamily},
            ymin=0, ymax=0.8,
            ytick={0,0.1,0.2,0.3,0.4,0.5},
            width=6.5cm, height=5.2cm,
            enlarge x limits=0.22,
            nodes near coords,
            nodes near coords align={vertical},
            every node near coord/.append style={font=\scriptsize},
          ]
            \addplot[fill=CellGreen, draw=DarkGreen] coordinates {
              (dormiu,0.00)(pulou,0.50)(caiu,0.33)(desceu,0.17)
            };
            \node[above, font=\scriptsize\bfseries, DarkGreen, align=center]
              at (axis cs:pulou, 0.60) {amostrado\\$\checkmark$};
          \end{axis}
        \end{tikzpicture}}
        \vspace{0.2em}
        \begin{block}{Garantia de correção}
          $p'$ concentra massa onde $M_p$ discorda de $M_q$, garantindo distribuição \textbf{idêntica a amostrar de $M_p$}.
        \end{block}
      \end{column}
    \end{columns}
  \end{frame}

  \begin{frame}{Decodificação Especulativa: Variantes e Resultados}
    \begin{columns}[T]
      \begin{column}{0.50\textwidth}
        \textbf{Leviathan et al.\ (2023)}:
        \begin{itemize}
          \item Testado: T5-XXL (11B) + T5-small (60M) como draft
          \item Speedup $2$--$3\times$ dependendo da tarefa
          \item HumanEval (código): maior speedup (tokens repetitivos)
          \item Sem degradação de qualidade
        \end{itemize}
      \end{column}
      \begin{column}{0.46\textwidth}
        \textbf{Variantes modernas}:
        \begin{itemize}
          \item \textbf{Self-speculative}: camadas rascunho do próprio modelo
          \item \textbf{Medusa}: múltiplas cabeças de predição em paralelo
          \item \textbf{Eagle}: alinha representações de draft e target
        \end{itemize}
        \vspace{0.3em}
        \begin{alertblock}{Requisito}
          Benefício depende da taxa de aceitação $\alpha$: quanto mais similares os modelos, maior $\alpha$ e maior o speedup.
        \end{alertblock}
      \end{column}
    \end{columns}
  \end{frame}

  %======================================================================================
  % SECTION 8: BENCHMARKS
  %======================================================================================
  \section{Benchmarks e Avaliação de LLMs}

  \begin{frame}{Perplexidade: A Métrica Mais Simples}
    \textbf{Definição}: mede o quão ``surpreso'' o modelo fica ao ver o texto.
    \[
      \mathrm{PPL}(w_1,\ldots,w_N)
        = \exp\!\left(-\frac{1}{N}\sum_{i=1}^{N}\log P(w_i \mid w_1,\ldots,w_{i-1})\right)
    \]
    \begin{itemize}
      \item Menor PPL $\to$ modelo estima melhor a distribuição do corpus
      \item PPL = 1: certeza absoluta; PPL = $|V|$: uniforme sobre o vocabulário
      \item Diretamente ligada à \textit{cross-entropy loss} do treino
    \end{itemize}
    \vspace{0.5em}
    \begin{exampleblock}{Intuição}
      PPL = 20 $\Rightarrow$ o modelo hesita, em média, entre 20 tokens igualmente prováveis a cada passo.
    \end{exampleblock}
  \end{frame}

  \begin{frame}{Perplexidade: Por que Não é Suficiente?}
    \begin{alertblock}{Limitações da perplexidade como única métrica}
      \begin{itemize}
        \item \textbf{Acurácia factual}: um modelo com PPL baixa pode ainda assim alucinar fatos
        \item \textbf{Raciocínio}: baixa surpresa no texto não implica capacidade de resolver problemas
        \item \textbf{Seguimento de instruções}: não captura se o modelo obedece comandos do usuário
        \item \textbf{Desempenho em tarefas reais}: PPL não prediz bem resultados em downstream tasks
      \end{itemize}
    \end{alertblock}
    \vspace{0.5em}
    \textcolor{DarkRed}{$\Rightarrow$ Precisamos de benchmarks específicos por tarefa para avaliar LLMs adequadamente.}
  \end{frame}

  \begin{frame}{Desafios na Avaliação de LLMs}
    \begin{columns}[T]
      \begin{column}{0.50\textwidth}
        Por que avaliar LLMs é difícil?
        \begin{itemize}
          \item Saída livre $\to$ métricas automáticas são imperfeitas
          \item \textit{Data contamination}: benchmarks no pré-treino
          \item Capacidades emergentes: falha em instâncias simples, acerta difíceis
          \item Benchmarks ficam saturados rapidamente
        \end{itemize}
      \end{column}
      \begin{column}{0.46\textwidth}
        \begin{block}{Categorias de avaliação}
          \begin{itemize}
            \item \textbf{Múltipla escolha}: ARC, MMLU
            \item \textbf{Geração + verificação}: GSM8k
            \item \textbf{Raciocínio}: HellaSwag, WinoGrande
            \item \textbf{Instrução}: MT-Bench, AlpacaEval
            \item \textbf{Arena humana}: Chatbot Arena (Elo)
          \end{itemize}
        \end{block}
      \end{column}
    \end{columns}
  \end{frame}

  \begin{frame}{ARC: AI2 Reasoning Challenge \cite{clark2018arc}}
    \begin{columns}[T]
      \begin{column}{0.52\textwidth}
        \textbf{Clark et al.\ (2018)}, Allen Institute for AI\\[0.4em]
        7787 questões de múltipla escolha de ciências (ensino fundamental/médio).\\[0.4em]
        \textbf{Duas partições}:
        \begin{itemize}
          \item \textbf{ARC-Easy}: respondidas por sistemas de recuperação simples
          \item \textbf{ARC-Challenge}: exigem raciocínio; falham em sistemas simples
        \end{itemize}
      \end{column}
      \begin{column}{0.44\textwidth}
        \textbf{Exemplo (ARC-Challenge)}:\\[0.2em]
        \textit{``Qual propriedade de um asteróide seria afetada de forma diferente pela gravidade de Júpiter e da Terra?''}\\
        (A) cor~~ (B) massa~~ (C) volume~~ \textbf{(D) peso}\\[0.5em]
        Avaliação: 0-shot ou 25-shot.\\[0.3em]
        \begin{exampleblock}{Resultados recentes}
          Llama 3 8B: 79.7\% (ARC-C 25-shot)\\
          GPT-4: $\approx$96\%
        \end{exampleblock}
      \end{column}
    \end{columns}
  \end{frame}

  \begin{frame}{GSM8k: Raciocínio Matemático \cite{cobbe2021gsm8k}}
    \begin{columns}[T]
      \begin{column}{0.52\textwidth}
        \textbf{Cobbe et al.\ (2021)}, OpenAI\\[0.4em]
        8500 problemas matemáticos de nível fundamental.\\[0.4em]
        \textbf{Características}:
        \begin{itemize}
          \item Resolução em múltiplos passos (2--8 operações)
          \item Resposta final é um número inteiro
          \item Avaliado com \textit{chain-of-thought} (CoT){\footnotesize, o modelo mostra o raciocínio passo a passo antes da resposta final}
        \end{itemize}
        Exemplo:\\
        \textit{``Natalia vendeu 48 clips em abril e metade disso em maio. Total?''}
        $\Rightarrow 48 + 24 = 72$
      \end{column}
      \begin{column}{0.44\textwidth}
        \resizebox{.8\linewidth}{!}{%
        \begin{tabular}{lr}
          \toprule
          \textbf{Modelo} & \textbf{GSM8k (\%)} \\
          \midrule
          GPT-2 (1.5B)    & $\approx$0   \\
          GPT-3 (175B)    & 17.9         \\
          Llama 2 7B      & 14.6         \\
          Llama 2 70B     & 56.8         \\
          Mixtral 8×7B    & 74.4         \\
          Llama 3 8B      & 79.6         \\
          Llama 3 70B     & 93.0         \\
          GPT-4           & 92.0         \\
          \bottomrule
        \end{tabular}}
      \end{column}
    \end{columns}
  \end{frame}

  \begin{frame}{MMLU: Massive Multitask Language Understanding \cite{hendrycks2021mmlu}}
    \begin{columns}[T]
      \begin{column}{0.50\textwidth}
        \textbf{Hendrycks et al.\ (2021)}, UC Berkeley\\[0.4em]
        15\,908 questões em \textbf{57 áreas}:
        \begin{itemize}
          \item Matemática, física, química, biologia
          \item Direito, medicina, economia
          \item História, filosofia, ciência da computação
        \end{itemize}
        Nível: colegial ao profissional avançado.\\[0.2em]
        Avaliação: 5-shot.\\[0.3em]
        \begin{exampleblock}{Exemplo: Conhecimento Clínico}
          {\footnotesize Qual achado laboratorial é mais típico de anemia ferropriva?\\[0.2em]
          (A) MCV elevado \quad (B) Ferritina elevada\\
          \textbf{(C) MCV reduzido e RDW elevado} \quad (D) B12 reduzida}
        \end{exampleblock}
      \end{column}
      \begin{column}{0.46\textwidth}
        \resizebox{.8\linewidth}{!}{%
        \begin{tabular}{lr}
          \toprule
          \textbf{Modelo} & \textbf{MMLU (\%)} \\
          \midrule
          GPT-3 (175B)       & 43.9 \\
          Llama 2 7B         & 45.3 \\
          Llama 2 70B        & 68.9 \\
          Mixtral 8×7B       & 70.6 \\
          Llama 3 8B         & 66.6 \\
          Llama 3 70B        & 79.5 \\
          GPT-4              & 86.4 \\
          \midrule
          Humano (estimado)  & $\approx$89.8 \\
          \bottomrule
        \end{tabular}}
      \end{column}
    \end{columns}
  \end{frame}

  \begin{frame}{HellaSwag e WinoGrande: Raciocínio de Senso Comum}
    \begin{columns}[T]
      \begin{column}{0.50\textwidth}
        \begin{block}{HellaSwag \cite{zellers2019hellaswag}}
          \textbf{Zellers et al.\ (2019)}\\
          Completar a continuação mais plausível.\\[0.2em]
          \textit{``Alguém lava um carro. Ele esfrega o capô\ldots''}\\
          (A) e joga fora.\\
          (B) por muito tempo.\\
          \textbf{(C) com esponja verde.}\\
          (D) e come um bolo.\\[0.3em]
          GPT-3: 78.9\%; Llama 3 8B: 82.0\%
        \end{block}
      \end{column}
      \begin{column}{0.46\textwidth}
        \begin{block}{WinoGrande}
          \textbf{Sakaguchi et al.\ (2021)}\\
          Resolução de co-referência adversarial.\\[0.2em]
          \textit{``Mark é mais alto que Paul, então \underline{ele} agachou na foto.''}\\
          (Mark ou Paul agachou?)\\[0.3em]
          Requer raciocínio físico/social.\\[0.3em]
          Llama 3 8B: 73.0\%; GPT-4: $\approx$87\%
        \end{block}
      \end{column}
    \end{columns}
  \end{frame}

  \begin{frame}{Open LLM Leaderboard: Como Comparar Modelos}
    \begin{columns}[T]
      \begin{column}{0.52\textwidth}
        Hugging Face mantém \textit{leaderboard} público de LLMs avaliados em condições padronizadas:
        \begin{itemize}
          \item ARC-Challenge (25-shot)
          \item HellaSwag (10-shot)
          \item MMLU (5-shot)
          \item TruthfulQA (0-shot)
          \item WinoGrande (5-shot)
          \item GSM8k (5-shot)
        \end{itemize}
        Média dos 6 = \textbf{score composto}.
      \end{column}
      \begin{column}{0.44\textwidth}
        \begin{alertblock}{Cuidados na interpretação}
          \begin{itemize}
            \item Modelos treinados para benchmarks
            \item Contaminação de dados de teste
            \item Prompts e few-shot afetam resultados
            \item Saturação rápida (MMLU-Pro, ARC-AGI)
          \end{itemize}
        \end{alertblock}
      \end{column}
    \end{columns}
  \end{frame}

  \begin{frame}{Comparativo de Modelos nos Principais Benchmarks}
    \begin{center}
    \resizebox{\linewidth}{!}{%
    \begin{tabular}{lrrrrrrr}
      \toprule
      \textbf{Modelo} & \textbf{ARC-C} & \textbf{HSwag} & \textbf{MMLU} & \textbf{TruthQA} & \textbf{WinoG} & \textbf{GSM8k} & \textbf{Média} \\
      \midrule
      Llama 2 7B   & 53.1 & 78.6 & 45.3 & 39.4 & 74.0 & 14.6 & 50.8 \\
      Llama 2 13B  & 59.4 & 82.1 & 54.8 & 37.4 & 77.0 & 28.7 & 56.6 \\
      Llama 2 70B  & 67.3 & 87.3 & 68.9 & 44.9 & 83.7 & 56.8 & 68.2 \\
      Mistral 7B   & 60.0 & 81.3 & 60.1 & 42.2 & 75.3 & 52.2 & 61.8 \\
      Mixtral 8×7B & 66.4 & 86.7 & 70.6 & 46.8 & 81.2 & 74.4 & 71.0 \\
      Llama 3 8B   & 79.7 & 82.0 & 66.6 & 44.0 & 73.0 & 79.6 & 70.8 \\
      Llama 3 70B  & 93.0 & 92.2 & 79.5 & 52.8 & 85.3 & 93.0 & 82.6 \\
      \bottomrule
    \end{tabular}}
    \end{center}
    \vspace{0.2em}
    \begin{center}
      \footnotesize Valores aproximados do Open LLM Leaderboard (2024). HSwag = HellaSwag; TruthQA = TruthfulQA; WinoG = WinoGrande.
    \end{center}
  \end{frame}

  \begin{frame}{Benchmarks Específicos por Tarefa}
    \begin{columns}[T]
      \begin{column}{0.50\textwidth}
        {\footnotesize Benchmarks gerais (MMLU, ARC\ldots) medem conhecimento amplo, mas \textbf{não capturam} capacidades especializadas:}
        \begin{block}{Contexto longo}
          {\footnotesize\begin{itemize}
            \item \textbf{RULER}: tarefas sintéticas em janelas de 4K--128K tokens
            \item \textbf{NIAH}: ``agulha'' escondida em textos de 100K+ tokens
          \end{itemize}}
        \end{block}
        \begin{block}{Chamada de ferramentas}
          {\footnotesize\begin{itemize}
            \item \textbf{BFCL}: 2000 pares função/argumento em múltiplas linguagens
            \item \textbf{ToolBench}: uso de APIs reais via ReAct
          \end{itemize}}
        \end{block}
      \end{column}
      \begin{column}{0.46\textwidth}
        \begin{block}{Código}
          {\footnotesize\begin{itemize}
            \item \textbf{HumanEval}: geração de código Python; métrica \textit{pass@k}
            \item \textbf{SWE-bench}: resolve issues reais do GitHub
          \end{itemize}}
        \end{block}
        \begin{block}{Outros domínios}
          {\footnotesize\begin{itemize}
            \item \textbf{MT-Bench}: instrução multi-turno avaliada por GPT-4
            \item \textbf{MATH}: 12500 problemas de competição
            \item \textbf{MMMU}: multimodal (imagem + texto), 57 disciplinas
          \end{itemize}}
        \end{block}
        \vspace{0.2em}
        {\footnotesize \textcolor{DarkBlue}{$\Rightarrow$ Nenhum benchmark único é suficiente: avaliação requer um portfólio de tarefas.}}
      \end{column}
    \end{columns}
  \end{frame}

  %======================================================================================
  % RESUMO
  %======================================================================================
  \section*{Resumo}

  \begin{frame}{Resumo da Aula}
    \resizebox{\linewidth}{!}{%
    \begin{tabular}{llll}
      \toprule
      \textbf{Tópico} & \textbf{Conceito-chave} & \textbf{Inovação / Referência} & \textbf{Modelo exemplo} \\
      \midrule
      Decoder-only      & Atenção causal, sem encoder   & GPT \cite{radford2018gpt1}        & GPT, Llama \\
      Leis de escala    & $20N$ tokens; $N^*\!\propto C^{0.5}$ & Chinchilla \cite{hoffmann2022chinchilla} & Llama 3 \\
      Normalização      & RMSNorm + Pre-Norm            & \cite{ba2016layernorm,zhang2019rmsnorm} & Llama 2/3 \\
      MHA $\to$ GQA     & Redução de KV-cache           & \cite{ainslie2023gqa}             & Llama 2 70B \\
      Sliding Window    & Atenção local $\mathcal{O}(nW)$ & \cite{jiang2023mistral}           & Mistral 7B \\
      MoE               & Top-$k$ routing, especialistas & \cite{shazeer2017moe,jiang2024mixtral} & Mixtral 8×7B \\
      Decod.\ especulativa & Draft + verify             & \cite{leviathan2023speculative}   & Llama 3 \\
      Benchmarks        & ARC, GSM8k, MMLU, HellaSwag   & \cite{clark2018arc,cobbe2021gsm8k,hendrycks2021mmlu} & --- \\
      \bottomrule
    \end{tabular}}
  \end{frame}

  %======================================================================================
  % REFERÊNCIAS
  %======================================================================================
  \begin{frame}[allowframebreaks]{Referências}
    \sloppy
    \printbibliography[heading=none]
  \end{frame}

\end{document}
